% Recovered from the previously completed PDF; no new mathematical content.
% Prose is editable; formula boxes retain original symbol positions.
% Build from repository root: xelatex -no-shell-escape -halt-on-error -output-directory=... print/fudan/NAME.tex
\documentclass[11pt,a4paper]{article}
\usepackage[margin=23mm,headheight=14pt,headsep=8.5mm,footskip=10mm]{geometry}
\usepackage{amsmath,amssymb}
\DeclareMathAlphabet{\mathsf}{OT1}{cmss}{m}{n}
\usepackage[no-math]{fontspec}
\usepackage{xeCJK}
\setmainfont{Latin Modern Roman}
\setCJKmainfont{FandolSong-Regular.otf}[Path=assets/fonts/fandol/,BoldFont=FandolSong-Bold.otf,ItalicFont=FandolSong-Regular.otf,Script=CustomDefault]
\xeCJKsetup{PunctStyle=quanjiao,CheckSingle=true}
\usepackage{setspace}\setstretch{1.12}
\setlength{\parindent}{2em}\setlength{\parskip}{3pt}
\emergencystretch=6em
\hfuzz=9pt
\usepackage{titlesec}
\titleformat{\section}{\centering\normalfont\Large\bfseries}{}{0pt}{}
\titlespacing*{\section}{0pt}{1.8\baselineskip}{1.0\baselineskip}
\usepackage{fancyhdr}\pagestyle{fancy}\fancyhf{}
\fancyhead[L]{\small 数学手记·高等代数}\fancyhead[R]{\small 22A05–22A11}
\fancyfoot[C]{\thepage}\renewcommand{\headrulewidth}{0.4pt}
\usepackage[hidelinks,unicode]{hyperref}
\hypersetup{pdftitle={22A05–22A11：题目、解答与小结},pdfauthor={Vincent Arqain}}
\newcommand{\RB}[4]{\raisebox{0pt}[#2bp][#3bp]{\hbox to #1bp{#4\hss}}}
\newcommand{\RG}[5]{\hbox to 0pt{\kern#1bp\raisebox{#2bp}[0pt][0pt]{\fontsize{#3bp}{#3bp}\csname font#4\endcsname\char#5}\hss}}
\newcommand{\RT}[4]{\hbox to 0pt{\kern#1bp\raisebox{#2bp}[0pt][0pt]{\fontsize{#3bp}{#3bp}\selectfont#4}\hss}}
\newcommand{\RS}[4]{\hbox to 0pt{\kern#1bp\raisebox{#2bp}[0pt][0pt]{\fontsize{#3bp}{#3bp}\selectfont\ensuremath{#4}}\hss}}
\newcommand{\RR}[4]{\hbox to 0pt{\kern#1bp\raisebox{#2bp}[0pt][0pt]{\rule{#3bp}{#4bp}}\hss}}
\newcommand{\fontmi}{\usefont{OML}{cmm}{m}{it}}
\newcommand{\fontsy}{\usefont{OMS}{cmsy}{m}{n}}
\newcommand{\fontex}{\usefont{OMX}{cmex}{m}{n}}
\newcommand{\fontrm}{\usefont{OT1}{cmr}{m}{n}}
\newcommand{\fontbf}{\usefont{OT1}{cmr}{bx}{n}}
\newcommand{\fontam}{\usefont{U}{msa}{m}{n}}
\newcommand{\fontbb}{\usefont{U}{msb}{m}{n}}
\widowpenalty=10000\clubpenalty=10000\raggedbottom
\begin{document}
\begin{center}{\LARGE\bfseries 22A05–22A11\par}\vspace{8pt}复旦大学高等代数习题集·题目与解答\end{center}
% Original page 1.
\section*{2022A05}
设\allowbreak{}\hskip.1em \RB{8.366}{8.310}{0.501}{%
\RG{0.000}{-0.001}{10.9091}{mi}{65}%
}\hskip.1em\allowbreak{} 为二阶实方阵。证明：存在二阶非异实方阵\allowbreak{}\hskip.1em \RB{8.573}{8.190}{2.617}{%
\RG{0.000}{-0.001}{10.9091}{mi}{81}%
}\hskip.1em\allowbreak{} ，使得

\[
\RB{63.322}{11.428}{2.616}{%
\RT{8.625}{4.504}{7.9701}{T}\RT{54.442}{4.504}{7.9701}{T}\RG{0.000}{-0.000}{10.9091}{mi}{81}\RG{14.909}{-0.000}{10.9091}{mi}{65}\RG{23.091}{-0.000}{10.9091}{mi}{81}\RG{34.745}{-0.000}{10.9091}{rm}{61}\RG{46.261}{-0.000}{10.9091}{mi}{65}\RG{60.728}{-0.000}{10.9091}{mi}{58}%
}
\]
证明：事实上，可以把\allowbreak{}\hskip.1em \RB{8.573}{8.192}{2.615}{%
\RG{0.000}{0.001}{10.9091}{mi}{81}%
}\hskip.1em\allowbreak{} 取为正交阵。将\allowbreak{}\hskip.1em \RB{8.366}{8.312}{0.500}{%
\RG{0.000}{0.001}{10.9091}{mi}{65}%
}\hskip.1em\allowbreak{} 分解为对称部分与反对称部分：

\[
\RB{231.947}{18.268}{7.978}{%
\RT{132.242}{11.344}{7.9701}{T}\RT{221.873}{11.344}{7.9701}{T}\RG{102.545}{7.385}{10.9091}{mi}{65}\RG{113.151}{7.385}{10.9091}{rm}{43}\RG{124.060}{7.385}{10.9091}{mi}{65}\RG{192.176}{7.385}{10.9091}{mi}{65}\RG{202.782}{7.385}{10.9091}{sy}{0}\RG{213.691}{7.385}{10.9091}{mi}{65}\RG{0.000}{0.005}{10.9091}{mi}{65}\RG{11.213}{0.005}{10.9091}{rm}{61}\RG{22.728}{0.005}{10.9091}{mi}{83}\RG{32.470}{0.005}{10.9091}{rm}{43}\RG{43.379}{0.005}{10.9091}{mi}{75}\RG{52.805}{0.005}{10.9091}{mi}{59}\RG{79.467}{0.005}{10.9091}{mi}{83}\RG{89.835}{0.005}{10.9091}{rm}{61}\RG{139.722}{0.005}{10.9091}{mi}{59}\RG{166.373}{0.005}{10.9091}{mi}{75}\RG{179.464}{0.005}{10.9091}{rm}{61}\RG{229.353}{0.005}{10.9091}{mi}{58}\RG{117.809}{-7.478}{10.9091}{rm}{50}\RG{207.439}{-7.478}{10.9091}{rm}{50}\RR{102.545}{2.732}{35.982}{0.436}\RR{192.176}{2.732}{35.981}{0.436}%
}
\]
\noindent \RB{391.770}{31.145}{14.128}{%
\RT{0.000}{21.917}{10.9091}{写}\RT{10.909}{21.917}{10.9091}{成}\RG{181.184}{18.652}{10.9092}{ex}{32}\RG{220.758}{18.652}{10.9092}{ex}{33}\RG{340.607}{18.652}{10.9092}{ex}{32}\RG{378.600}{18.652}{10.9092}{ex}{33}\RG{189.820}{7.277}{10.9091}{mi}{117}\RG{210.547}{7.277}{10.9091}{mi}{118}\RG{353.486}{7.277}{10.9091}{rm}{48}\RG{373.144}{7.277}{10.9091}{rm}{49}\RG{72.801}{-0.003}{10.9091}{mi}{83}\RG{83.149}{-0.003}{10.9091}{rm}{61}\RG{94.664}{-0.003}{10.9091}{mi}{11}\RG{101.667}{-0.003}{10.9091}{mi}{73}\RG{109.759}{-0.003}{10.9091}{rm}{43}\RG{120.668}{-0.003}{10.9091}{mi}{72}\RG{130.006}{-0.003}{10.9091}{mi}{59}\RG{156.668}{-0.003}{10.9091}{mi}{72}\RG{169.668}{-0.003}{10.9091}{rm}{61}\RG{231.212}{-0.003}{10.9091}{mi}{59}\RG{257.863}{-0.003}{10.9091}{mi}{75}\RG{270.955}{-0.003}{10.9091}{rm}{61}\RG{282.470}{-0.003}{10.9091}{mi}{116}\RG{286.409}{-0.003}{10.9091}{mi}{74}\RG{292.300}{-0.003}{10.9091}{mi}{59}\RG{318.940}{-0.003}{10.9091}{mi}{74}\RG{329.092}{-0.003}{10.9091}{rm}{61}\RG{389.055}{-0.003}{10.9091}{mi}{59}\RG{190.103}{-7.898}{10.9091}{mi}{118}\RG{206.028}{-7.898}{10.9091}{sy}{0}\RG{214.513}{-7.898}{10.9091}{mi}{117}\RG{349.244}{-7.898}{10.9091}{sy}{0}\RG{357.728}{-7.898}{10.9091}{rm}{49}\RG{373.143}{-7.898}{10.9091}{rm}{48}%
}

其中\allowbreak{}\hskip.1em \RB{58.931}{7.975}{2.603}{%
\RG{0.000}{0.002}{10.9091}{mi}{11}\RG{7.004}{0.002}{10.9091}{mi}{59}\RG{11.847}{0.002}{10.9091}{mi}{117}\RG{18.087}{0.002}{10.9091}{mi}{59}\RG{22.931}{0.002}{10.9091}{mi}{118}\RG{28.593}{0.002}{10.9091}{mi}{59}\RG{33.437}{0.002}{10.9091}{mi}{116}\RG{40.458}{0.002}{10.9091}{sy}{50}\RG{50.761}{0.002}{10.9091}{bb}{82}%
}\hskip.1em\allowbreak{} 。直接相乘可得

\[
\RB{163.866}{10.307}{3.224}{%
\RG{9.954}{4.507}{7.9701}{rm}{50}\RG{39.720}{4.507}{7.9701}{rm}{50}\RG{63.464}{4.507}{7.9701}{rm}{50}\RG{0.000}{0.003}{10.9091}{mi}{72}\RG{17.717}{0.003}{10.9091}{rm}{61}\RG{29.215}{0.003}{10.9091}{rm}{40}\RG{33.475}{0.003}{10.9091}{mi}{117}\RG{46.876}{0.003}{10.9091}{rm}{43}\RG{57.785}{0.003}{10.9091}{mi}{118}\RG{68.198}{0.003}{10.9091}{rm}{41}\RG{72.440}{0.003}{10.9091}{mi}{73}\RG{78.080}{0.003}{10.9091}{mi}{59}\RG{104.742}{0.003}{10.9091}{mi}{72}\RG{114.691}{0.003}{10.9091}{mi}{74}\RG{124.842}{0.003}{10.9091}{rm}{61}\RG{136.357}{0.003}{10.9091}{sy}{0}\RG{144.842}{0.003}{10.9091}{mi}{74}\RG{151.933}{0.003}{10.9091}{mi}{72}\RG{161.271}{0.003}{10.9091}{mi}{58}%
}
\]
当\allowbreak{}\hskip.1em \RB{63.456}{8.686}{3.223}{%
\RG{0.000}{0.004}{10.9091}{rm}{40}\RG{4.242}{0.004}{10.9091}{mi}{117}\RG{10.482}{0.004}{10.9091}{mi}{59}\RG{15.315}{0.004}{10.9091}{mi}{118}\RG{21.015}{0.004}{10.9091}{rm}{41}\RG{28.288}{0.004}{10.9091}{sy}{54}\RG{28.288}{0.004}{10.9091}{rm}{61}\RG{39.787}{0.004}{10.9091}{rm}{40}\RG{44.019}{0.004}{10.9091}{rm}{48}\RG{49.500}{0.004}{10.9091}{mi}{59}\RG{54.348}{0.004}{10.9091}{rm}{48}\RG{59.803}{0.004}{10.9091}{rm}{41}%
}\hskip.1em\allowbreak{} 时，令

\[
\RB{123.816}{19.591}{14.126}{%
\RG{71.192}{18.654}{10.9092}{ex}{32}\RG{110.766}{18.654}{10.9092}{ex}{33}\RG{43.544}{7.379}{10.9091}{rm}{49}\RG{79.829}{7.279}{10.9091}{mi}{117}\RG{100.556}{7.279}{10.9091}{mi}{118}\RG{24.365}{0.327}{10.9091}{sy}{112}\RG{0.000}{-0.001}{10.9091}{mi}{81}\RG{11.654}{-0.001}{10.9091}{rm}{61}\RG{121.221}{-0.001}{10.9091}{mi}{58}\RG{39.701}{-5.720}{7.9701}{rm}{50}\RG{63.446}{-5.720}{7.9701}{rm}{50}\RG{80.112}{-7.896}{10.9091}{mi}{118}\RG{96.036}{-7.896}{10.9091}{sy}{0}\RG{104.521}{-7.896}{10.9091}{mi}{117}\RG{33.456}{-8.872}{10.9091}{mi}{117}\RG{46.857}{-8.872}{10.9091}{rm}{43}\RG{57.766}{-8.872}{10.9091}{mi}{118}\RR{24.365}{2.726}{43.813}{0.436}\RR{33.456}{0.545}{34.722}{0.436}%
}
\]
于是\allowbreak{}\hskip.1em \RB{8.573}{8.190}{2.617}{%
\RG{0.000}{-0.001}{10.9091}{mi}{81}%
}\hskip.1em\allowbreak{} T\allowbreak{}\hskip.1em \RB{20.088}{8.190}{2.617}{%
\RG{0.000}{-0.001}{10.9091}{rm}{61}\RG{11.515}{-0.001}{10.9091}{mi}{81}%
}\hskip.1em\allowbreak{} 、\allowbreak{}\hskip.1em \RB{33.824}{9.758}{2.617}{%
\RG{0.000}{-0.001}{10.9091}{mi}{81}\RG{8.624}{3.958}{7.9701}{rm}{50}\RG{16.387}{-0.001}{10.9091}{rm}{61}\RG{27.902}{-0.001}{10.9091}{mi}{73}%
}\hskip.1em\allowbreak{} ，所以\allowbreak{}\hskip.1em \RB{8.573}{8.190}{2.617}{%
\RG{0.000}{-0.001}{10.9091}{mi}{81}%
}\hskip.1em\allowbreak{} 为实正交阵。又因为\allowbreak{}\hskip.1em \RB{57.978}{8.190}{1.406}{%
\RG{0.000}{-0.001}{10.9091}{mi}{83}\RG{10.348}{-0.001}{10.9091}{rm}{61}\RG{21.864}{-0.001}{10.9091}{mi}{11}\RG{28.868}{-0.001}{10.9091}{mi}{73}\RG{36.958}{-0.001}{10.9091}{rm}{43}\RG{47.867}{-0.001}{10.9091}{mi}{72}%
}\hskip.1em\allowbreak{} 、\allowbreak{}\hskip.1em \RB{8.573}{8.190}{2.617}{%
\RG{0.000}{-0.001}{10.9091}{mi}{81}%
}\hskip.1em\allowbreak{} 是\allowbreak{}\hskip.1em \RB{10.111}{7.950}{0.501}{%
\RG{0.000}{-0.001}{10.9091}{mi}{72}%
}\hskip.1em\allowbreak{} 的实数倍，故

\[
\RB{136.047}{8.191}{2.616}{%
\RG{0.000}{-0.000}{10.9091}{mi}{81}\RG{8.618}{-0.000}{10.9091}{mi}{83}\RG{18.973}{-0.000}{10.9091}{rm}{61}\RG{30.488}{-0.000}{10.9091}{mi}{83}\RG{37.797}{-0.000}{10.9091}{mi}{81}\RG{46.416}{-0.000}{10.9091}{mi}{59}\RG{73.078}{-0.000}{10.9091}{mi}{81}\RG{81.696}{-0.000}{10.9091}{mi}{75}\RG{94.798}{-0.000}{10.9091}{rm}{61}\RG{106.313}{-0.000}{10.9091}{sy}{0}\RG{114.798}{-0.000}{10.9091}{mi}{75}\RG{124.834}{-0.000}{10.9091}{mi}{81}\RG{133.453}{-0.000}{10.9091}{mi}{58}%
}
\]
从而

\[
\RB{108.564}{11.431}{37.065}{%
\RT{8.625}{4.507}{7.9701}{T}\RG{0.000}{0.003}{10.9091}{mi}{81}\RG{14.909}{0.003}{10.9091}{mi}{65}\RG{23.091}{0.003}{10.9091}{mi}{81}\RG{34.745}{0.003}{10.9091}{rm}{61}\RG{46.260}{0.003}{10.9091}{mi}{81}\RG{54.885}{0.003}{10.9091}{rm}{40}\RG{59.127}{0.003}{10.9091}{mi}{83}\RG{68.870}{0.003}{10.9091}{rm}{43}\RG{79.780}{0.003}{10.9091}{mi}{75}\RG{89.825}{0.003}{10.9091}{rm}{41}\RG{94.067}{0.003}{10.9091}{mi}{81}\RG{62.204}{-13.657}{7.9701}{rm}{50}\RG{98.939}{-13.657}{7.9701}{rm}{50}\RG{34.745}{-18.161}{10.9091}{rm}{61}\RG{46.260}{-18.161}{10.9091}{mi}{83}\RG{53.570}{-18.161}{10.9091}{mi}{81}\RG{69.360}{-18.161}{10.9091}{sy}{0}\RG{80.269}{-18.161}{10.9091}{mi}{75}\RG{90.305}{-18.161}{10.9091}{mi}{81}\RT{99.686}{-31.821}{7.9701}{T}\RG{34.745}{-36.325}{10.9091}{rm}{61}\RG{46.260}{-36.325}{10.9091}{mi}{83}\RG{56.003}{-36.325}{10.9091}{sy}{0}\RG{66.913}{-36.325}{10.9091}{mi}{75}\RG{79.989}{-36.325}{10.9091}{rm}{61}\RG{91.504}{-36.325}{10.9091}{mi}{65}\RG{105.970}{-36.325}{10.9091}{mi}{58}%
}
\]
当\allowbreak{}\hskip.1em \RB{46.517}{7.761}{0.744}{%
\RG{0.000}{-0.004}{10.9091}{mi}{117}\RG{9.275}{-0.004}{10.9091}{rm}{61}\RG{20.790}{-0.004}{10.9091}{mi}{118}\RG{29.500}{-0.004}{10.9091}{rm}{61}\RG{40.999}{-0.004}{10.9091}{rm}{48}%
}\hskip.1em\allowbreak{} 时，\allowbreak{}\hskip.1em \RB{34.790}{8.187}{0.744}{%
\RG{0.000}{-0.004}{10.9091}{mi}{83}\RG{10.349}{-0.004}{10.9091}{rm}{61}\RG{21.864}{-0.004}{10.9091}{mi}{11}\RG{28.868}{-0.004}{10.9091}{mi}{73}%
}\hskip.1em\allowbreak{} 。取

\[
\RB{74.332}{19.594}{14.122}{%
\RG{23.170}{18.658}{10.9092}{ex}{32}\RG{61.163}{18.658}{10.9092}{ex}{33}\RG{31.806}{7.283}{10.9091}{rm}{49}\RG{51.465}{7.283}{10.9091}{rm}{48}\RG{0.000}{0.003}{10.9091}{mi}{81}\RG{11.655}{0.003}{10.9091}{rm}{61}\RG{71.618}{0.003}{10.9091}{mi}{59}\RG{31.806}{-7.892}{10.9091}{rm}{48}\RG{47.223}{-7.892}{10.9091}{sy}{0}\RG{55.708}{-7.892}{10.9091}{rm}{49}%
}
\]
则仍有\allowbreak{}\hskip.1em \RB{8.573}{8.194}{2.613}{%
\RG{0.000}{0.003}{10.9091}{mi}{81}%
}\hskip.1em\allowbreak{} T\allowbreak{}\hskip.1em \RB{29.091}{8.194}{2.613}{%
\RG{0.000}{0.003}{10.9091}{mi}{81}\RG{11.654}{0.003}{10.9091}{rm}{61}\RG{23.169}{0.003}{10.9091}{mi}{73}%
}\hskip.1em\allowbreak{} 、\allowbreak{}\hskip.1em \RB{8.573}{8.194}{2.613}{%
\RG{0.000}{0.003}{10.9091}{mi}{81}%
}\hskip.1em\allowbreak{} T\allowbreak{}\hskip.1em \RB{38.025}{8.194}{2.613}{%
\RG{0.000}{0.003}{10.9091}{mi}{83}\RG{7.309}{0.003}{10.9091}{mi}{81}\RG{18.973}{0.003}{10.9091}{rm}{61}\RG{30.488}{0.003}{10.9091}{mi}{83}%
}\hskip.1em\allowbreak{} 、\allowbreak{}\hskip.1em \RB{8.573}{8.194}{2.613}{%
\RG{0.000}{0.003}{10.9091}{mi}{81}%
}\hskip.1em\allowbreak{} T\allowbreak{}\hskip.1em \RB{46.159}{8.194}{2.613}{%
\RG{0.000}{0.003}{10.9091}{mi}{74}\RG{7.091}{0.003}{10.9091}{mi}{81}\RG{18.753}{0.003}{10.9091}{rm}{61}\RG{30.268}{0.003}{10.9091}{sy}{0}\RG{38.753}{0.003}{10.9091}{mi}{74}%
}\hskip.1em\allowbreak{} ，因而同样得到\allowbreak{}\hskip.1em \RB{8.573}{8.194}{2.613}{%
\RG{0.000}{0.003}{10.9091}{mi}{81}%
}\hskip.1em\allowbreak{} T\allowbreak{}\hskip.1em \RB{84.640}{8.314}{2.613}{%
\RG{0.000}{0.003}{10.9091}{mi}{65}\RG{8.182}{0.003}{10.9091}{mi}{81}\RG{19.837}{0.003}{10.9091}{rm}{61}\RG{31.352}{0.003}{10.9091}{mi}{83}\RG{40.934}{0.003}{10.9091}{sy}{0}\RG{51.683}{0.003}{10.9091}{mi}{75}\RG{64.760}{0.003}{10.9091}{rm}{61}\RG{76.275}{0.003}{10.9091}{mi}{65}%
}\hskip.1em\allowbreak{} T。两种情形涵盖所有二阶实方阵。\allowbreak{}\hskip.1em \RB{8.366}{7.993}{0.502}{%
\RG{0.000}{-0.002}{10.9091}{am}{3}%
}\hskip.1em\allowbreak{}

方法小结：转置保持对称部分，并使反对称部分变号。构造一个正交反射，使它保持前者、翻转后者，就得到所需的合同变换；上述证明还给出了\allowbreak{}\hskip.1em \RB{8.573}{8.192}{2.615}{%
\RG{0.000}{0.001}{10.9091}{mi}{81}%
}\hskip.1em\allowbreak{} 的显式公式。

题目来源：谢启鸿《复旦大学高等代数习题集（2023版）》，第124页。

% Original page 2.
\section*{2022A06}
设\allowbreak{}\hskip.1em \RB{21.762}{8.310}{2.606}{%
\RG{0.000}{-0.001}{10.9091}{mi}{65}\RG{8.182}{-0.001}{10.9091}{mi}{59}\RG{13.015}{-0.001}{10.9091}{mi}{66}%
}\hskip.1em\allowbreak{} 为\allowbreak{}\hskip.1em \RB{6.729}{5.321}{0.621}{%
\RG{0.000}{-0.001}{10.9091}{mi}{110}%
}\hskip.1em\allowbreak{} 阶实对称阵，满足\allowbreak{}\hskip.1em \RB{61.643}{9.758}{0.501}{%
\RG{0.000}{-0.001}{10.9091}{mi}{65}\RG{8.181}{3.958}{7.9701}{rm}{50}\RG{12.914}{-0.001}{10.9091}{mi}{66}\RG{24.766}{-0.001}{10.9091}{rm}{61}\RG{36.281}{-0.001}{10.9091}{mi}{65}\RG{44.463}{-0.001}{10.9091}{mi}{66}\RG{53.278}{-0.001}{10.9091}{mi}{65}%
}\hskip.1em\allowbreak{} 。证明：\allowbreak{}\hskip.1em \RB{48.729}{8.310}{0.501}{%
\RG{0.000}{-0.001}{10.9091}{mi}{65}\RG{8.182}{-0.001}{10.9091}{mi}{66}\RG{20.034}{-0.001}{10.9091}{rm}{61}\RG{31.549}{-0.001}{10.9091}{mi}{66}\RG{40.364}{-0.001}{10.9091}{mi}{65}%
}\hskip.1em\allowbreak{} 。

证明：令\allowbreak{}\hskip.1em \RB{70.641}{8.307}{0.744}{%
\RG{0.000}{-0.004}{10.9091}{mi}{67}\RG{11.607}{-0.004}{10.9091}{rm}{61}\RG{23.123}{-0.004}{10.9091}{mi}{65}\RG{31.305}{-0.004}{10.9091}{mi}{66}\RG{42.551}{-0.004}{10.9091}{sy}{0}\RG{53.460}{-0.004}{10.9091}{mi}{66}\RG{62.275}{-0.004}{10.9091}{mi}{65}%
}\hskip.1em\allowbreak{} 。我们把结论化为\allowbreak{}\hskip.1em \RB{8.791}{8.187}{0.744}{%
\RG{0.000}{-0.004}{10.9091}{mi}{67}%
}\hskip.1em\allowbreak{} 的全部元素平方之和为零。

由\allowbreak{}\hskip.1em \RB{61.643}{9.760}{0.500}{%
\RG{0.000}{0.001}{10.9091}{mi}{65}\RG{8.182}{3.960}{7.9701}{rm}{50}\RG{12.914}{0.001}{10.9091}{mi}{66}\RG{24.766}{0.001}{10.9091}{rm}{61}\RG{36.281}{0.001}{10.9091}{mi}{65}\RG{44.463}{0.001}{10.9091}{mi}{66}\RG{53.278}{0.001}{10.9091}{mi}{65}%
}\hskip.1em\allowbreak{} ，在右边乘以\allowbreak{}\hskip.1em \RB{8.747}{7.952}{0.500}{%
\RG{0.000}{0.001}{10.9091}{mi}{66}%
}\hskip.1em\allowbreak{} ，得

\[
\RB{77.722}{10.306}{2.603}{%
\RG{8.182}{4.506}{7.9701}{rm}{50}\RG{21.736}{4.506}{7.9701}{rm}{50}\RG{0.000}{0.002}{10.9091}{mi}{65}\RG{12.914}{0.002}{10.9091}{mi}{66}\RG{29.499}{0.002}{10.9091}{rm}{61}\RG{41.014}{0.002}{10.9091}{mi}{65}\RG{49.196}{0.002}{10.9091}{mi}{66}\RG{58.011}{0.002}{10.9091}{mi}{65}\RG{66.193}{0.002}{10.9091}{mi}{66}\RG{75.007}{0.002}{10.9091}{mi}{59}%
}
\]
因而

\[
\RB{111.626}{10.302}{3.229}{%
\RG{20.944}{4.502}{7.9701}{rm}{50}\RG{34.499}{4.502}{7.9701}{rm}{50}\RT{0.000}{-0.002}{10.9091}{t}\RT{4.244}{-0.002}{10.9091}{r}\RG{8.520}{-0.002}{10.9091}{rm}{40}\RG{12.763}{-0.002}{10.9091}{mi}{65}\RG{25.677}{-0.002}{10.9091}{mi}{66}\RG{39.231}{-0.002}{10.9091}{rm}{41}\RG{46.486}{-0.002}{10.9091}{rm}{61}\RT{58.019}{-0.002}{10.9091}{t}\RT{62.263}{-0.002}{10.9091}{r}\RG{66.540}{-0.002}{10.9091}{rm}{40}\RG{70.782}{-0.002}{10.9091}{mi}{65}\RG{78.963}{-0.002}{10.9091}{mi}{66}\RG{87.778}{-0.002}{10.9091}{mi}{65}\RG{95.960}{-0.002}{10.9091}{mi}{66}\RG{104.790}{-0.002}{10.9091}{rm}{41}\RG{109.032}{-0.002}{10.9091}{mi}{58}%
}
\]
由于\allowbreak{}\hskip.1em \RB{8.366}{8.311}{0.500}{%
\RG{0.000}{0.000}{10.9091}{mi}{65}%
}\hskip.1em\allowbreak{} T\allowbreak{}\hskip.1em \RB{19.881}{8.311}{0.500}{%
\RG{0.000}{0.000}{10.9091}{rm}{61}\RG{11.515}{0.000}{10.9091}{mi}{65}%
}\hskip.1em\allowbreak{} 、\allowbreak{}\hskip.1em \RB{8.747}{7.951}{0.500}{%
\RG{0.000}{0.000}{10.9091}{mi}{66}%
}\hskip.1em\allowbreak{} T\allowbreak{}\hskip.1em \RB{20.262}{7.951}{0.500}{%
\RG{0.000}{0.000}{10.9091}{rm}{61}\RG{11.515}{0.000}{10.9091}{mi}{66}%
}\hskip.1em\allowbreak{} ，有\allowbreak{}\hskip.1em \RB{8.791}{8.191}{0.740}{%
\RG{0.000}{0.000}{10.9091}{mi}{67}%
}\hskip.1em\allowbreak{} T\allowbreak{}\hskip.1em \RB{58.782}{8.311}{0.500}{%
\RG{0.000}{0.000}{10.9091}{rm}{61}\RG{11.515}{0.000}{10.9091}{mi}{66}\RG{20.330}{0.000}{10.9091}{mi}{65}\RG{30.943}{0.000}{10.9091}{sy}{0}\RG{41.852}{0.000}{10.9091}{mi}{65}\RG{50.034}{0.000}{10.9091}{mi}{66}%
}\hskip.1em\allowbreak{} 。展开可得

\[
\RB{290.912}{29.590}{37.070}{%
\RG{63.510}{26.998}{10.9092}{ex}{0}\RG{180.162}{26.998}{10.9092}{ex}{1}\RT{21.340}{22.666}{7.9701}{T}\RT{0.000}{18.161}{10.9091}{t}\RT{4.244}{18.161}{10.9091}{r}\RG{8.520}{18.161}{10.9091}{rm}{40}\RG{12.762}{18.161}{10.9091}{mi}{67}\RG{27.624}{18.161}{10.9091}{mi}{67}\RG{36.201}{18.161}{10.9091}{rm}{41}\RG{43.456}{18.161}{10.9091}{rm}{61}\RT{54.989}{18.161}{10.9091}{t}\RT{59.233}{18.161}{10.9091}{r}\RG{68.510}{18.161}{10.9091}{rm}{40}\RG{72.753}{18.161}{10.9091}{mi}{66}\RG{81.567}{18.161}{10.9091}{mi}{65}\RG{92.181}{18.161}{10.9091}{sy}{0}\RG{103.090}{18.161}{10.9091}{mi}{65}\RG{111.272}{18.161}{10.9091}{mi}{66}\RG{120.094}{18.161}{10.9091}{rm}{41}\RG{124.327}{18.161}{10.9091}{rm}{40}\RG{128.578}{18.161}{10.9091}{mi}{65}\RG{136.760}{18.161}{10.9091}{mi}{66}\RG{148.006}{18.161}{10.9091}{sy}{0}\RG{158.915}{18.161}{10.9091}{mi}{66}\RG{167.730}{18.161}{10.9091}{mi}{65}\RG{175.919}{18.161}{10.9091}{rm}{41}\RG{84.756}{4.502}{7.9701}{rm}{50}\RG{274.345}{4.502}{7.9701}{rm}{50}\RG{43.474}{-0.002}{10.9091}{rm}{61}\RT{54.989}{-0.002}{10.9091}{t}\RT{59.233}{-0.002}{10.9091}{r}\RG{63.510}{-0.002}{10.9091}{rm}{40}\RG{67.753}{-0.002}{10.9091}{mi}{66}\RG{76.567}{-0.002}{10.9091}{mi}{65}\RG{89.489}{-0.002}{10.9091}{mi}{66}\RG{98.311}{-0.002}{10.9091}{rm}{41}\RG{104.977}{-0.002}{10.9091}{sy}{0}\RT{115.886}{-0.002}{10.9091}{t}\RT{120.130}{-0.002}{10.9091}{r}\RG{124.406}{-0.002}{10.9091}{rm}{40}\RG{128.649}{-0.002}{10.9091}{mi}{66}\RG{137.464}{-0.002}{10.9091}{mi}{65}\RG{145.646}{-0.002}{10.9091}{mi}{66}\RG{154.460}{-0.002}{10.9091}{mi}{65}\RG{162.656}{-0.002}{10.9091}{rm}{41}\RG{169.323}{-0.002}{10.9091}{sy}{0}\RT{180.232}{-0.002}{10.9091}{t}\RT{184.476}{-0.002}{10.9091}{r}\RG{188.753}{-0.002}{10.9091}{rm}{40}\RG{192.996}{-0.002}{10.9091}{mi}{65}\RG{201.177}{-0.002}{10.9091}{mi}{66}\RG{209.992}{-0.002}{10.9091}{mi}{65}\RG{218.174}{-0.002}{10.9091}{mi}{66}\RG{227.003}{-0.002}{10.9091}{rm}{41}\RG{233.657}{-0.002}{10.9091}{rm}{43}\RT{244.579}{-0.002}{10.9091}{t}\RT{248.823}{-0.002}{10.9091}{r}\RG{253.099}{-0.002}{10.9091}{rm}{40}\RG{257.341}{-0.002}{10.9091}{mi}{65}\RG{265.523}{-0.002}{10.9091}{mi}{66}\RG{279.078}{-0.002}{10.9091}{mi}{65}\RG{287.259}{-0.002}{10.9091}{rm}{41}\RG{83.207}{-13.662}{7.9701}{rm}{50}\RG{96.761}{-13.662}{7.9701}{rm}{50}\RG{43.474}{-18.166}{10.9091}{rm}{61}\RG{54.973}{-18.166}{10.9091}{rm}{50}\RT{62.262}{-18.166}{10.9091}{t}\RT{66.506}{-18.166}{10.9091}{r}\RG{70.783}{-18.166}{10.9091}{rm}{40}\RG{75.025}{-18.166}{10.9091}{mi}{65}\RG{87.939}{-18.166}{10.9091}{mi}{66}\RG{101.494}{-18.166}{10.9091}{rm}{41}\RG{108.160}{-18.166}{10.9091}{sy}{0}\RG{119.069}{-18.166}{10.9091}{rm}{50}\RT{126.342}{-18.166}{10.9091}{t}\RT{130.586}{-18.166}{10.9091}{r}\RG{134.862}{-18.166}{10.9091}{rm}{40}\RG{139.105}{-18.166}{10.9091}{mi}{65}\RG{147.287}{-18.166}{10.9091}{mi}{66}\RG{156.102}{-18.166}{10.9091}{mi}{65}\RG{164.284}{-18.166}{10.9091}{mi}{66}\RG{173.112}{-18.166}{10.9091}{rm}{41}\RG{43.474}{-36.330}{10.9091}{rm}{61}\RG{54.973}{-36.330}{10.9091}{rm}{48}\RG{60.444}{-36.330}{10.9091}{mi}{58}%
}
\]
这里使用了迹的循环性质。其依据是：对任意同阶方阵\allowbreak{}\hskip.1em \RB{22.929}{7.952}{2.604}{%
\RG{0.000}{0.001}{10.9091}{mi}{88}\RG{9.273}{0.001}{10.9091}{mi}{59}\RG{14.116}{0.001}{10.9091}{mi}{89}%
}\hskip.1em\allowbreak{} ，

\[
\RB{204.948}{10.866}{15.211}{%
\RG{50.201}{10.366}{10.9092}{ex}{88}\RG{108.437}{10.366}{10.9092}{ex}{88}\RT{0.000}{0.002}{10.9091}{t}\RT{4.244}{0.002}{10.9091}{r}\RG{8.520}{0.002}{10.9091}{rm}{40}\RG{12.762}{0.002}{10.9091}{mi}{88}\RG{22.646}{0.002}{10.9091}{mi}{89}\RG{31.414}{0.002}{10.9091}{rm}{41}\RG{38.669}{0.002}{10.9091}{rm}{61}\RG{67.778}{0.002}{10.9091}{mi}{120}\RG{81.278}{0.002}{10.9091}{mi}{121}\RG{96.922}{0.002}{10.9091}{rm}{61}\RG{126.013}{0.002}{10.9091}{mi}{121}\RG{138.627}{0.002}{10.9091}{mi}{120}\RG{155.157}{0.002}{10.9091}{rm}{61}\RT{166.673}{0.002}{10.9091}{t}\RT{170.917}{0.002}{10.9091}{r}\RG{175.193}{0.002}{10.9091}{rm}{40}\RG{179.435}{0.002}{10.9091}{mi}{89}\RG{188.185}{0.002}{10.9091}{mi}{88}\RG{198.088}{0.002}{10.9091}{rm}{41}\RG{202.320}{0.002}{10.9091}{rm}{59}\RG{74.013}{-1.634}{7.9701}{mi}{105}\RG{76.890}{-1.634}{7.9701}{mi}{106}\RG{86.627}{-1.634}{7.9701}{mi}{106}\RG{90.500}{-1.634}{7.9701}{mi}{105}\RG{131.362}{-1.634}{7.9701}{mi}{106}\RG{135.236}{-1.634}{7.9701}{mi}{105}\RG{144.862}{-1.634}{7.9701}{mi}{105}\RG{147.739}{-1.634}{7.9701}{mi}{106}\RG{53.520}{-13.085}{7.9701}{mi}{105}\RG{56.398}{-13.085}{7.9701}{mi}{59}\RG{58.749}{-13.085}{7.9701}{mi}{106}\RG{111.992}{-13.085}{7.9701}{mi}{106}\RG{115.403}{-13.085}{7.9701}{mi}{59}\RG{117.754}{-13.085}{7.9701}{mi}{105}%
}
\]
把多个相邻因子合并为一项，即可循环移动乘积中的因子。

最后，若\allowbreak{}\hskip.1em \RB{43.004}{8.684}{3.761}{%
\RG{0.000}{0.002}{10.9091}{mi}{67}\RG{11.608}{0.002}{10.9091}{rm}{61}\RG{23.106}{0.002}{10.9091}{rm}{40}\RG{27.365}{0.002}{10.9091}{mi}{99}\RG{32.086}{-1.635}{7.9701}{mi}{105}\RG{34.963}{-1.635}{7.9701}{mi}{106}\RG{39.351}{0.002}{10.9091}{rm}{41}%
}\hskip.1em\allowbreak{} ，则

\[
\RB{104.722}{17.652}{15.213}{%
\RG{60.299}{13.637}{7.9701}{mi}{110}\RG{77.876}{13.637}{7.9701}{mi}{110}\RG{54.989}{10.364}{10.9092}{ex}{88}\RG{72.566}{10.364}{10.9092}{ex}{88}\RT{21.340}{4.505}{7.9701}{T}\RG{94.863}{4.505}{7.9701}{rm}{50}\RT{0.000}{0.001}{10.9091}{t}\RT{4.244}{0.001}{10.9091}{r}\RG{8.520}{0.001}{10.9091}{rm}{40}\RG{12.763}{0.001}{10.9091}{mi}{67}\RG{27.624}{0.001}{10.9091}{mi}{67}\RG{36.202}{0.001}{10.9091}{rm}{41}\RG{43.457}{0.001}{10.9091}{rm}{61}\RG{90.142}{0.001}{10.9091}{mi}{99}\RG{102.128}{0.001}{10.9091}{mi}{58}\RG{94.863}{-2.696}{7.9701}{mi}{105}\RG{97.740}{-2.696}{7.9701}{mi}{106}\RG{56.016}{-13.087}{7.9701}{mi}{105}\RG{58.899}{-13.087}{7.9701}{rm}{61}\RG{65.483}{-13.087}{7.9701}{rm}{49}\RG{73.093}{-13.087}{7.9701}{mi}{106}\RG{76.977}{-13.087}{7.9701}{rm}{61}\RG{83.560}{-13.087}{7.9701}{rm}{49}%
}
\]
各项均为非负实数，其和为零迫使每个\allowbreak{}\hskip.1em \RB{32.034}{7.762}{3.766}{%
\RG{0.000}{-0.004}{10.9091}{mi}{99}\RG{4.721}{-1.640}{7.9701}{mi}{105}\RG{7.598}{-1.640}{7.9701}{mi}{106}\RG{15.017}{-0.004}{10.9091}{rm}{61}\RG{26.515}{-0.004}{10.9091}{rm}{48}%
}\hskip.1em\allowbreak{} 。于是\allowbreak{}\hskip.1em \RB{28.625}{8.187}{0.744}{%
\RG{0.000}{-0.004}{10.9091}{mi}{67}\RG{11.608}{-0.004}{10.9091}{rm}{61}\RG{23.106}{-0.004}{10.9091}{rm}{48}%
}\hskip.1em\allowbreak{} ，即\allowbreak{}\hskip.1em \RB{48.730}{8.307}{0.504}{%
\RG{0.000}{-0.004}{10.9091}{mi}{65}\RG{8.182}{-0.004}{10.9091}{mi}{66}\RG{20.034}{-0.004}{10.9091}{rm}{61}\RG{31.550}{-0.004}{10.9091}{mi}{66}\RG{40.365}{-0.004}{10.9091}{mi}{65}%
}\hskip.1em\allowbreak{} 。\allowbreak{}\hskip.1em \RB{8.366}{7.991}{0.504}{%
\RG{0.000}{-0.004}{10.9091}{am}{3}%
}\hskip.1em\allowbreak{}

方法小结：遇到要证两个实对称阵交换的问题，可以考察交换子\allowbreak{}\hskip.1em \RB{47.933}{8.315}{0.500}{%
\RG{0.000}{0.004}{10.9091}{mi}{65}\RG{8.182}{0.004}{10.9091}{mi}{66}\RG{19.636}{0.004}{10.9091}{sy}{0}\RG{30.753}{0.004}{10.9091}{mi}{66}\RG{39.568}{0.004}{10.9091}{mi}{65}%
}\hskip.1em\allowbreak{} 的Frobenius范数平方。条件先给出迹的相等关系，再由“平方和为零”推出矩阵为零。

题目来源：谢启鸿《复旦大学高等代数习题集（2023版）》，第124页。

% Original page 3.
\section*{2022A07}
固定分块方式\allowbreak{}\hskip.1em \RB{83.065}{6.859}{2.358}{%
\RG{0.000}{-0.001}{10.9091}{mi}{110}\RG{9.578}{-0.001}{10.9091}{rm}{61}\RG{21.094}{-0.001}{10.9091}{mi}{110}\RG{27.642}{-1.638}{7.9701}{rm}{49}\RG{34.798}{-0.001}{10.9091}{rm}{43}\RG{45.707}{-0.001}{10.9091}{sy}{1}\RG{50.540}{-0.001}{10.9091}{sy}{1}\RG{55.384}{-0.001}{10.9091}{sy}{1}\RG{60.860}{-0.001}{10.9091}{rm}{43}\RG{71.769}{-0.001}{10.9091}{mi}{110}\RG{78.317}{-1.778}{7.9701}{mi}{107}%
}\hskip.1em\allowbreak{} 。设\allowbreak{}\hskip.1em \RB{7.504}{8.190}{0.741}{%
\RG{0.000}{-0.001}{10.9091}{sy}{83}%
}\hskip.1em\allowbreak{} 为数域\allowbreak{}\hskip.1em \RB{8.889}{7.972}{0.501}{%
\RG{0.000}{-0.001}{10.9091}{bb}{75}%
}\hskip.1em\allowbreak{} 上所有形如

\[
\RB{135.787}{11.253}{55.182}{%
\RG{22.727}{10.327}{10.9092}{ex}{48}\RG{128.927}{10.327}{10.9092}{ex}{49}\RG{32.273}{0.004}{10.9091}{mi}{65}\RG{59.384}{0.004}{10.9091}{mi}{65}\RG{86.496}{0.004}{10.9091}{sy}{1}\RG{91.329}{0.004}{10.9091}{sy}{1}\RG{96.172}{0.004}{10.9091}{sy}{1}\RG{111.198}{0.004}{10.9091}{mi}{65}\RG{40.455}{-1.632}{7.9701}{rm}{49}\RG{44.688}{-1.632}{7.9701}{rm}{49}\RG{67.567}{-1.632}{7.9701}{rm}{49}\RG{71.799}{-1.632}{7.9701}{rm}{50}\RG{119.379}{-1.773}{7.9701}{rm}{49}\RG{123.614}{-1.773}{7.9701}{mi}{107}\RG{22.727}{-8.873}{10.9092}{ex}{66}\RG{128.927}{-8.873}{10.9092}{ex}{67}\RG{38.120}{-15.171}{10.9091}{rm}{48}\RG{59.384}{-15.171}{10.9091}{mi}{65}\RG{86.496}{-15.171}{10.9091}{sy}{1}\RG{91.328}{-15.171}{10.9091}{sy}{1}\RG{96.172}{-15.171}{10.9091}{sy}{1}\RG{111.198}{-15.171}{10.9091}{mi}{65}\RG{22.727}{-15.419}{10.9092}{ex}{66}\RG{128.927}{-15.419}{10.9092}{ex}{67}\RG{67.567}{-16.807}{7.9701}{rm}{50}\RG{71.799}{-16.807}{7.9701}{rm}{50}\RG{119.379}{-16.948}{7.9701}{rm}{50}\RG{123.614}{-16.948}{7.9701}{mi}{107}\RG{22.727}{-21.964}{10.9092}{ex}{66}\RG{128.927}{-21.964}{10.9092}{ex}{67}\RG{0.000}{-24.692}{10.9091}{mi}{65}\RG{11.212}{-24.692}{10.9091}{rm}{61}\RT{39.331}{-26.858}{10.9091}{.}\RT{118.449}{-26.858}{10.9091}{.}\RT{62.198}{-27.854}{10.9091}{.}\RT{88.007}{-27.854}{10.9091}{.}\RG{22.727}{-28.510}{10.9092}{ex}{66}\RG{128.927}{-28.510}{10.9092}{ex}{67}\RT{39.331}{-30.843}{10.9091}{.}\RT{66.443}{-30.843}{10.9091}{.}\RT{92.252}{-30.843}{10.9091}{.}\RT{118.449}{-30.843}{10.9091}{.}\RT{70.687}{-33.831}{10.9091}{.}\RT{96.497}{-33.831}{10.9091}{.}\RT{39.331}{-34.828}{10.9091}{.}\RT{118.449}{-34.828}{10.9091}{.}\RG{22.727}{-35.492}{10.9092}{ex}{64}\RG{128.927}{-35.492}{10.9092}{ex}{65}\RG{38.120}{-50.003}{10.9091}{rm}{48}\RG{60.685}{-50.003}{10.9091}{sy}{1}\RG{65.518}{-50.003}{10.9091}{sy}{1}\RG{70.362}{-50.003}{10.9091}{sy}{1}\RG{91.041}{-50.003}{10.9091}{rm}{48}\RG{111.004}{-50.003}{10.9091}{mi}{65}\RG{119.186}{-51.780}{7.9701}{mi}{107}\RG{123.801}{-51.780}{7.9701}{mi}{107}%
}
\]
的分块上三角阵构成的集合，其中\allowbreak{}\hskip.1em \RB{14.850}{8.311}{3.762}{%
\RG{0.000}{-0.000}{10.9091}{mi}{65}\RG{8.181}{-1.636}{7.9701}{mi}{105}\RG{11.058}{-1.636}{7.9701}{mi}{106}%
}\hskip.1em\allowbreak{} 为\allowbreak{}\hskip.1em \RB{33.307}{5.857}{3.762}{%
\RG{0.000}{-0.000}{10.9091}{mi}{110}\RG{6.548}{-1.636}{7.9701}{mi}{105}\RG{12.206}{0.001}{10.9091}{sy}{2}\RG{22.967}{-0.000}{10.9091}{mi}{110}\RG{29.515}{-1.636}{7.9701}{mi}{106}%
}\hskip.1em\allowbreak{} 矩阵。证明：加法、数乘、乘法、求逆（对可逆矩阵而言）和求伴随均保持这一分块上三角形状。这里“伴随”指古典伴随矩阵adj\allowbreak{}\hskip.1em \RB{16.077}{8.687}{3.222}{%
\RG{0.000}{0.005}{10.9091}{rm}{40}\RG{4.242}{0.005}{10.9091}{mi}{65}\RG{12.424}{0.005}{10.9091}{rm}{41}%
}\hskip.1em\allowbreak{} 。

证明：令\allowbreak{}\hskip.1em \RB{88.707}{7.954}{3.759}{%
\RG{0.000}{0.003}{10.9091}{mi}{78}\RG{8.765}{-1.633}{7.9701}{mi}{106}\RG{16.177}{0.003}{10.9091}{rm}{61}\RG{27.693}{0.003}{10.9091}{mi}{110}\RG{34.241}{-1.633}{7.9701}{rm}{49}\RG{41.397}{0.003}{10.9091}{rm}{43}\RG{52.306}{0.003}{10.9091}{sy}{1}\RG{57.139}{0.003}{10.9091}{sy}{1}\RG{61.983}{0.003}{10.9091}{sy}{1}\RG{67.459}{0.003}{10.9091}{rm}{43}\RG{78.368}{0.003}{10.9091}{mi}{110}\RG{84.916}{-1.633}{7.9701}{mi}{106}%
}\hskip.1em\allowbreak{} ，并设

\[
\RB{219.530}{9.297}{4.620}{%
\RG{135.310}{4.508}{7.9701}{mi}{110}\RG{0.000}{0.004}{10.9091}{mi}{87}\RG{17.715}{0.004}{10.9091}{rm}{61}\RT{29.230}{0.004}{10.9091}{s}\RT{33.528}{0.004}{10.9091}{p}\RT{39.594}{0.004}{10.9091}{a}\RT{45.049}{0.004}{10.9091}{n}\RG{51.114}{0.004}{10.9091}{sy}{102}\RG{56.568}{0.004}{10.9091}{mi}{101}\RG{66.381}{0.004}{10.9091}{mi}{59}\RG{71.213}{0.004}{10.9091}{mi}{58}\RG{76.057}{0.004}{10.9091}{mi}{58}\RG{80.901}{0.004}{10.9091}{mi}{58}\RG{85.733}{0.004}{10.9091}{mi}{59}\RG{90.577}{0.004}{10.9091}{mi}{101}\RG{106.826}{0.004}{10.9091}{sy}{103}\RG{115.302}{0.004}{10.9091}{sy}{18}\RG{126.825}{0.004}{10.9091}{bb}{75}\RG{162.765}{0.004}{10.9091}{rm}{40}\RG{166.998}{0.004}{10.9091}{rm}{49}\RG{175.492}{0.004}{10.9091}{sy}{20}\RG{187.008}{0.004}{10.9091}{mi}{106}\RG{195.155}{0.004}{10.9091}{sy}{20}\RG{206.671}{0.004}{10.9091}{mi}{107}\RG{212.693}{0.004}{10.9091}{rm}{41}\RG{216.936}{0.004}{10.9091}{mi}{58}\RG{10.303}{-1.632}{7.9701}{mi}{106}\RG{61.649}{-1.632}{7.9701}{rm}{49}\RG{95.703}{-1.685}{7.9701}{mi}{78}\RG{102.425}{-2.900}{5.9776}{mi}{106}%
}
\]
矩阵\allowbreak{}\hskip.1em \RB{8.366}{8.316}{0.500}{%
\RG{0.000}{0.005}{10.9091}{mi}{65}%
}\hskip.1em\allowbreak{} 属于\allowbreak{}\hskip.1em \RB{7.504}{8.196}{0.735}{%
\RG{0.000}{0.005}{10.9091}{sy}{83}%
}\hskip.1em\allowbreak{} ，当且仅当\allowbreak{}\hskip.1em \RB{59.992}{8.687}{3.757}{%
\RG{0.000}{0.005}{10.9091}{mi}{65}\RG{8.181}{0.005}{10.9091}{rm}{40}\RG{12.424}{0.005}{10.9091}{mi}{87}\RG{22.727}{-1.631}{7.9701}{mi}{106}\RG{27.109}{0.005}{10.9091}{rm}{41}\RG{34.382}{0.005}{10.9091}{sy}{18}\RG{45.897}{0.005}{10.9091}{mi}{87}\RG{56.200}{-1.631}{7.9701}{mi}{106}%
}\hskip.1em\allowbreak{} 对每个\allowbreak{}\hskip.1em \RB{4.831}{7.716}{2.731}{%
\RG{0.000}{0.005}{10.9091}{mi}{106}%
}\hskip.1em\allowbreak{} 均成立：这一条件恰好要求各个对角分块以下的分块全为零。

若\allowbreak{}\hskip.1em \RB{42.673}{8.316}{2.601}{%
\RG{0.000}{0.005}{10.9091}{mi}{65}\RG{8.182}{0.005}{10.9091}{mi}{59}\RG{13.015}{0.005}{10.9091}{mi}{66}\RG{24.882}{0.005}{10.9091}{sy}{50}\RG{35.170}{0.005}{10.9091}{sy}{83}%
}\hskip.1em\allowbreak{} 、\allowbreak{}\hskip.1em \RB{26.943}{7.978}{0.931}{%
\RG{0.000}{0.005}{10.9091}{mi}{99}\RG{7.751}{0.005}{10.9091}{sy}{50}\RG{18.054}{0.005}{10.9091}{bb}{75}%
}\hskip.1em\allowbreak{} ，则对每个\allowbreak{}\hskip.1em \RB{4.831}{7.716}{2.731}{%
\RG{0.000}{0.005}{10.9091}{mi}{106}%
}\hskip.1em\allowbreak{} ，

\[
\RB{204.435}{8.678}{22.475}{%
\RG{0.000}{-0.004}{10.9091}{rm}{40}\RG{4.243}{-0.004}{10.9091}{mi}{65}\RG{14.849}{-0.004}{10.9091}{rm}{43}\RG{25.758}{-0.004}{10.9091}{mi}{66}\RG{34.580}{-0.004}{10.9091}{rm}{41}\RG{38.813}{-0.004}{10.9091}{rm}{40}\RG{43.065}{-0.004}{10.9091}{mi}{87}\RG{57.750}{-0.004}{10.9091}{rm}{41}\RG{65.023}{-0.004}{10.9091}{sy}{18}\RG{76.539}{-0.004}{10.9091}{mi}{87}\RG{91.224}{-0.004}{10.9091}{mi}{59}\RG{117.890}{-0.004}{10.9091}{rm}{40}\RG{122.133}{-0.004}{10.9091}{mi}{99}\RG{126.846}{-0.004}{10.9091}{mi}{65}\RG{135.036}{-0.004}{10.9091}{rm}{41}\RG{139.269}{-0.004}{10.9091}{rm}{40}\RG{143.520}{-0.004}{10.9091}{mi}{87}\RG{158.206}{-0.004}{10.9091}{rm}{41}\RG{165.478}{-0.004}{10.9091}{sy}{18}\RG{176.993}{-0.004}{10.9091}{mi}{87}\RG{191.680}{-0.004}{10.9091}{mi}{59}\RG{53.368}{-1.640}{7.9701}{mi}{106}\RG{86.842}{-1.640}{7.9701}{mi}{106}\RG{153.824}{-1.640}{7.9701}{mi}{106}\RG{187.298}{-1.640}{7.9701}{mi}{106}\RG{84.720}{-9.331}{10.9092}{ex}{0}\RG{121.712}{-9.331}{10.9092}{ex}{1}\RG{13.333}{-18.168}{10.9091}{rm}{40}\RG{17.576}{-18.168}{10.9091}{mi}{65}\RG{25.758}{-18.168}{10.9091}{mi}{66}\RG{34.580}{-18.168}{10.9091}{rm}{41}\RG{38.813}{-18.168}{10.9091}{rm}{40}\RG{43.065}{-18.168}{10.9091}{mi}{87}\RG{57.750}{-18.168}{10.9091}{rm}{41}\RG{65.005}{-18.168}{10.9091}{rm}{61}\RG{76.539}{-18.168}{10.9091}{mi}{65}\RG{89.720}{-18.168}{10.9091}{mi}{66}\RG{98.542}{-18.168}{10.9091}{rm}{40}\RG{102.785}{-18.168}{10.9091}{mi}{87}\RG{117.470}{-18.168}{10.9091}{rm}{41}\RG{129.743}{-18.168}{10.9091}{sy}{18}\RG{141.258}{-18.168}{10.9091}{mi}{65}\RG{149.440}{-18.168}{10.9091}{rm}{40}\RG{153.682}{-18.168}{10.9091}{mi}{87}\RG{168.367}{-18.168}{10.9091}{rm}{41}\RG{175.640}{-18.168}{10.9091}{sy}{18}\RG{187.156}{-18.168}{10.9091}{mi}{87}\RG{201.841}{-18.168}{10.9091}{mi}{58}\RG{53.368}{-19.804}{7.9701}{mi}{106}\RG{113.088}{-19.804}{7.9701}{mi}{106}\RG{163.985}{-19.804}{7.9701}{mi}{106}\RG{197.459}{-19.804}{7.9701}{mi}{106}%
}
\]
故\allowbreak{}\hskip.1em \RB{69.830}{8.314}{2.603}{%
\RG{0.000}{0.003}{10.9091}{mi}{65}\RG{10.606}{0.003}{10.9091}{rm}{43}\RG{21.515}{0.003}{10.9091}{mi}{66}\RG{30.330}{0.003}{10.9091}{mi}{59}\RG{35.163}{0.003}{10.9091}{mi}{99}\RG{39.875}{0.003}{10.9091}{mi}{65}\RG{48.057}{0.003}{10.9091}{mi}{59}\RG{52.901}{0.003}{10.9091}{mi}{65}\RG{61.083}{0.003}{10.9091}{mi}{66}%
}\hskip.1em\allowbreak{} 都属于\allowbreak{}\hskip.1em \RB{7.504}{8.194}{0.737}{%
\RG{0.000}{0.003}{10.9091}{sy}{83}%
}\hskip.1em\allowbreak{} 。

若\allowbreak{}\hskip.1em \RB{29.003}{8.309}{0.938}{%
\RG{0.000}{-0.002}{10.9091}{mi}{65}\RG{11.212}{-0.002}{10.9091}{sy}{50}\RG{21.499}{-0.002}{10.9091}{sy}{83}%
}\hskip.1em\allowbreak{} 可逆，则\allowbreak{}\hskip.1em \RB{8.366}{8.309}{0.502}{%
\RG{0.000}{-0.002}{10.9091}{mi}{65}%
}\hskip.1em\allowbreak{} 在\allowbreak{}\hskip.1em \RB{14.095}{7.949}{3.764}{%
\RG{0.000}{-0.002}{10.9091}{mi}{87}\RG{10.303}{-1.638}{7.9701}{mi}{106}%
}\hskip.1em\allowbreak{} 上的限制是单射，所以

\[
\RB{103.185}{9.292}{3.763}{%
\RT{0.000}{-0.001}{10.9091}{d}\RT{6.066}{-0.001}{10.9091}{i}\RT{9.098}{-0.001}{10.9091}{m}\RG{20.004}{-0.001}{10.9091}{mi}{65}\RG{28.185}{-0.001}{10.9091}{rm}{40}\RG{32.428}{-0.001}{10.9091}{mi}{87}\RG{47.113}{-0.001}{10.9091}{rm}{41}\RG{54.368}{-0.001}{10.9091}{rm}{61}\RT{65.902}{-0.001}{10.9091}{d}\RT{71.967}{-0.001}{10.9091}{i}\RT{75.000}{-0.001}{10.9091}{m}\RG{85.905}{-0.001}{10.9091}{mi}{87}\RG{100.591}{-0.001}{10.9091}{mi}{58}\RG{42.731}{-1.637}{7.9701}{mi}{106}\RG{96.208}{-1.637}{7.9701}{mi}{106}%
}
\]
结合\allowbreak{}\hskip.1em \RB{59.992}{8.682}{3.762}{%
\RG{0.000}{-0.000}{10.9091}{mi}{65}\RG{8.182}{-0.000}{10.9091}{rm}{40}\RG{12.424}{-0.000}{10.9091}{mi}{87}\RG{22.727}{-1.636}{7.9701}{mi}{106}\RG{27.109}{-0.000}{10.9091}{rm}{41}\RG{34.382}{-0.000}{10.9091}{sy}{18}\RG{45.897}{-0.000}{10.9091}{mi}{87}\RG{56.200}{-1.636}{7.9701}{mi}{106}%
}\hskip.1em\allowbreak{} ，得到\allowbreak{}\hskip.1em \RB{59.993}{8.682}{3.762}{%
\RG{0.000}{-0.000}{10.9091}{mi}{65}\RG{8.182}{-0.000}{10.9091}{rm}{40}\RG{12.425}{-0.000}{10.9091}{mi}{87}\RG{22.728}{-1.636}{7.9701}{mi}{106}\RG{27.110}{-0.000}{10.9091}{rm}{41}\RG{34.365}{-0.000}{10.9091}{rm}{61}\RG{45.898}{-0.000}{10.9091}{mi}{87}\RG{56.201}{-1.636}{7.9701}{mi}{106}%
}\hskip.1em\allowbreak{} ，于是\allowbreak{}\hskip.1em \RB{71.311}{9.759}{3.762}{%
\RG{0.000}{-0.000}{10.9091}{mi}{65}\RG{8.182}{3.959}{7.9701}{sy}{0}\RG{14.768}{3.959}{7.9701}{rm}{49}\RG{19.500}{-0.000}{10.9091}{rm}{40}\RG{23.743}{-0.000}{10.9091}{mi}{87}\RG{34.046}{-1.636}{7.9701}{mi}{106}\RG{38.428}{-0.000}{10.9091}{rm}{41}\RG{45.683}{-0.000}{10.9091}{rm}{61}\RG{57.216}{-0.000}{10.9091}{mi}{87}\RG{67.520}{-1.636}{7.9701}{mi}{106}%
}\hskip.1em\allowbreak{} 。因此\allowbreak{}\hskip.1em \RB{40.322}{9.759}{0.936}{%
\RG{0.000}{-0.000}{10.9091}{mi}{65}\RG{8.182}{3.959}{7.9701}{sy}{0}\RG{14.768}{3.959}{7.9701}{rm}{49}\RG{22.531}{-0.000}{10.9091}{sy}{50}\RG{32.819}{-0.000}{10.9091}{sy}{83}%
}\hskip.1em\allowbreak{} 。

下面处理伴随矩阵，包括\allowbreak{}\hskip.1em \RB{8.366}{8.316}{0.500}{%
\RG{0.000}{0.005}{10.9091}{mi}{65}%
}\hskip.1em\allowbreak{} 不可逆的情形。对任意\allowbreak{}\hskip.1em \RB{26.162}{7.978}{0.931}{%
\RG{0.000}{0.005}{10.9091}{mi}{116}\RG{6.970}{0.005}{10.9091}{sy}{50}\RG{17.273}{0.005}{10.9091}{bb}{75}%
}\hskip.1em\allowbreak{} ，有\allowbreak{}\hskip.1em \RB{50.612}{8.316}{1.400}{%
\RG{0.000}{0.005}{10.9091}{mi}{65}\RG{9.949}{0.005}{10.9091}{rm}{43}\RG{20.200}{0.005}{10.9091}{mi}{116}\RG{24.138}{0.005}{10.9091}{mi}{73}\RG{32.821}{0.005}{10.9091}{sy}{50}\RG{43.109}{0.005}{10.9091}{sy}{83}%
}\hskip.1em\allowbreak{} 。多项式det\allowbreak{}\hskip.1em \RB{37.686}{8.687}{3.222}{%
\RG{0.000}{0.005}{10.9091}{rm}{40}\RG{4.243}{0.005}{10.9091}{mi}{65}\RG{14.191}{0.005}{10.9091}{rm}{43}\RG{24.443}{0.005}{10.9091}{mi}{116}\RG{28.381}{0.005}{10.9091}{mi}{73}\RG{34.033}{0.005}{10.9091}{rm}{41}%
}\hskip.1em\allowbreak{}的最高次项为\allowbreak{}\hskip.1em \RB{9.318}{7.974}{0.620}{%
\RG{0.000}{-0.000}{10.9091}{mi}{116}\RG{3.940}{3.959}{7.9701}{mi}{110}%
}\hskip.1em\allowbreak{} ，所以只有有限个根。除去这些根，

\[
\RB{198.001}{10.306}{3.226}{%
\RG{163.330}{4.506}{7.9701}{sy}{0}\RG{169.917}{4.506}{7.9701}{rm}{49}\RT{0.000}{0.001}{10.9091}{a}\RT{5.455}{0.001}{10.9091}{d}\RT{11.520}{0.001}{10.9091}{j}\RG{14.858}{0.001}{10.9091}{rm}{40}\RG{19.101}{0.001}{10.9091}{mi}{65}\RG{29.707}{0.001}{10.9091}{rm}{43}\RG{40.616}{0.001}{10.9091}{mi}{116}\RG{44.555}{0.001}{10.9091}{mi}{73}\RG{50.207}{0.001}{10.9091}{rm}{41}\RG{57.462}{0.001}{10.9091}{rm}{61}\RT{68.996}{0.001}{10.9091}{d}\RT{75.061}{0.001}{10.9091}{e}\RT{79.905}{0.001}{10.9091}{t}\RG{84.148}{0.001}{10.9091}{rm}{40}\RG{88.391}{0.001}{10.9091}{mi}{65}\RG{98.997}{0.001}{10.9091}{rm}{43}\RG{109.906}{0.001}{10.9091}{mi}{116}\RG{113.844}{0.001}{10.9091}{mi}{73}\RG{119.497}{0.001}{10.9091}{rm}{41}\RG{123.730}{0.001}{10.9091}{rm}{40}\RG{127.982}{0.001}{10.9091}{mi}{65}\RG{138.588}{0.001}{10.9091}{rm}{43}\RG{149.497}{0.001}{10.9091}{mi}{116}\RG{153.435}{0.001}{10.9091}{mi}{73}\RG{159.088}{0.001}{10.9091}{rm}{41}\RG{177.679}{0.001}{10.9091}{sy}{50}\RG{187.967}{0.001}{10.9091}{sy}{83}\RG{195.407}{0.001}{10.9091}{mi}{58}%
}
\]
伴随矩阵的每个元素都是代数余子式，因而是关于\allowbreak{}\hskip.1em \RB{4.100}{7.331}{0.618}{%
\RG{0.000}{0.002}{10.9091}{mi}{116}%
}\hskip.1em\allowbreak{} 的多项式。其对角分块以下的每个元素，在无穷多个\allowbreak{}\hskip.1em \RB{26.162}{7.970}{0.939}{%
\RG{0.000}{-0.003}{10.9091}{mi}{116}\RG{6.970}{-0.003}{10.9091}{sy}{50}\RG{17.273}{-0.003}{10.9091}{bb}{75}%
}\hskip.1em\allowbreak{} 处都等于零，故该多项式恒为零。令\allowbreak{}\hskip.1em \RB{23.987}{7.763}{0.743}{%
\RG{0.000}{-0.003}{10.9091}{mi}{116}\RG{6.970}{-0.003}{10.9091}{rm}{61}\RG{18.468}{-0.003}{10.9091}{rm}{48}%
}\hskip.1em\allowbreak{} ，便得

\[
\RB{54.877}{9.291}{3.229}{%
\RT{0.000}{-0.002}{10.9091}{a}\RT{5.455}{-0.002}{10.9091}{d}\RT{11.520}{-0.002}{10.9091}{j}\RG{14.858}{-0.002}{10.9091}{rm}{40}\RG{19.100}{-0.002}{10.9091}{mi}{65}\RG{27.282}{-0.002}{10.9091}{rm}{41}\RG{34.555}{-0.002}{10.9091}{sy}{50}\RG{44.843}{-0.002}{10.9091}{sy}{83}\RG{52.282}{-0.002}{10.9091}{mi}{58}%
}
\]
这就证明了全部封闭性。\allowbreak{}\hskip.1em \RB{8.366}{7.994}{0.501}{%
\RG{0.000}{-0.001}{10.9091}{am}{3}%
}\hskip.1em\allowbreak{}

方法小结：分块上三角形状等价于保持一串嵌套子空间。可逆性将子空间的包含关系提升为相等；伴随矩阵在奇异情形下的结论则由多项式恒等原理补足。

题目来源：谢启鸿《复旦大学高等代数习题集（2023版）》，第124页。

% Original page 4.
\section*{2022A08}
设\allowbreak{}\hskip.1em \RB{143.935}{8.681}{3.228}{%
\RG{0.000}{-0.001}{10.9091}{mi}{86}\RG{11.818}{-0.001}{10.9091}{rm}{61}\RG{23.333}{-0.001}{10.9091}{sy}{102}\RG{28.788}{-0.001}{10.9091}{mi}{122}\RG{37.371}{-0.001}{10.9091}{sy}{50}\RG{47.674}{-0.001}{10.9091}{bb}{67}\RG{58.583}{-0.001}{10.9091}{rm}{58}\RG{64.644}{-0.001}{10.9091}{sy}{106}\RG{67.675}{-0.001}{10.9091}{mi}{122}\RG{73.228}{-0.001}{10.9091}{sy}{106}\RG{79.288}{-0.001}{10.9091}{rm}{61}\RG{90.786}{-0.001}{10.9091}{rm}{49}\RG{96.258}{-0.001}{10.9091}{mi}{59}\RG{104.723}{-0.001}{10.9091}{mi}{122}\RG{113.323}{-0.001}{10.9091}{sy}{54}\RG{113.323}{-0.001}{10.9091}{rm}{61}\RG{124.838}{-0.001}{10.9091}{sy}{0}\RG{133.323}{-0.001}{10.9091}{rm}{49}\RG{138.777}{-0.001}{10.9091}{sy}{103}%
}\hskip.1em\allowbreak{} ，在\allowbreak{}\hskip.1em \RB{8.889}{7.950}{0.741}{%
\RG{0.000}{-0.001}{10.9091}{mi}{86}%
}\hskip.1em\allowbreak{} 上定义

\[
\RB{348.683}{8.682}{18.091}{%
\RG{48.686}{-0.001}{10.9091}{sy}{0}\RG{57.171}{-0.001}{10.9091}{rm}{49}\RG{65.048}{-0.001}{10.9091}{rm}{43}\RG{75.959}{-0.001}{10.9091}{mi}{122}\RG{88.189}{-0.001}{10.9091}{rm}{43}\RG{99.098}{-0.001}{10.9091}{mi}{122}\RG{111.328}{-0.001}{10.9091}{rm}{43}\RG{122.226}{-0.001}{10.9091}{rm}{51}\RG{127.692}{-0.001}{10.9091}{mi}{122}\RG{137.497}{-0.001}{10.9091}{mi}{122}\RG{212.786}{-0.001}{10.9091}{rm}{40}\RG{217.018}{-0.001}{10.9091}{rm}{49}\RG{224.907}{-0.000}{10.9091}{sy}{0}\RG{235.816}{-0.001}{10.9091}{mi}{107}\RG{241.839}{-0.000}{10.9091}{rm}{41}\RG{248.494}{-0.000}{10.9091}{rm}{43}\RG{259.392}{-0.000}{10.9091}{rm}{40}\RG{263.625}{-0.000}{10.9091}{rm}{49}\RG{271.501}{-0.000}{10.9091}{rm}{43}\RG{282.445}{-0.001}{10.9091}{mi}{107}\RG{288.467}{-0.000}{10.9091}{rm}{41}\RG{292.710}{-0.001}{10.9091}{mi}{122}\RG{81.033}{-1.637}{7.9701}{rm}{49}\RG{104.172}{-1.637}{7.9701}{rm}{50}\RG{132.765}{-1.637}{7.9701}{rm}{49}\RG{142.571}{-1.637}{7.9701}{rm}{50}\RG{0.000}{-7.381}{10.9091}{mi}{122}\RG{12.230}{-7.381}{10.9091}{sy}{8}\RG{23.139}{-7.381}{10.9091}{mi}{122}\RG{35.975}{-7.381}{10.9091}{rm}{61}\RG{148.498}{-7.381}{10.9091}{mi}{59}\RG{175.149}{-7.381}{10.9091}{mi}{107}\RG{183.612}{-7.381}{10.9091}{sy}{14}\RG{191.492}{-7.381}{10.9091}{mi}{122}\RG{200.075}{-7.381}{10.9091}{rm}{61}\RG{310.368}{-7.381}{10.9091}{rm}{40}\RG{314.611}{-7.381}{10.9091}{mi}{107}\RG{323.664}{-7.381}{10.9091}{sy}{50}\RG{333.967}{-7.381}{10.9091}{bb}{82}\RG{341.846}{-7.381}{10.9091}{rm}{41}\RG{346.088}{-7.381}{10.9091}{mi}{58}\RG{5.074}{-9.017}{7.9701}{rm}{49}\RG{28.213}{-9.017}{7.9701}{rm}{50}\RG{55.655}{-14.864}{10.9091}{rm}{51}\RG{63.532}{-14.864}{10.9091}{rm}{43}\RG{74.444}{-14.864}{10.9091}{mi}{122}\RG{86.674}{-14.864}{10.9091}{rm}{43}\RG{97.583}{-14.864}{10.9091}{mi}{122}\RG{109.813}{-14.864}{10.9091}{sy}{0}\RG{120.722}{-14.864}{10.9091}{mi}{122}\RG{130.528}{-14.864}{10.9091}{mi}{122}\RG{212.786}{-14.864}{10.9091}{rm}{40}\RG{217.019}{-14.864}{10.9091}{rm}{49}\RG{224.895}{-14.864}{10.9091}{rm}{43}\RG{235.816}{-14.864}{10.9091}{mi}{107}\RG{241.839}{-14.864}{10.9091}{rm}{41}\RG{248.494}{-14.864}{10.9091}{rm}{43}\RG{259.392}{-14.864}{10.9091}{rm}{40}\RG{263.625}{-14.864}{10.9091}{rm}{49}\RG{271.535}{-14.864}{10.9091}{sy}{0}\RG{282.445}{-14.864}{10.9091}{mi}{107}\RG{288.467}{-14.864}{10.9091}{rm}{41}\RG{292.710}{-14.864}{10.9091}{mi}{122}\RG{79.518}{-16.500}{7.9701}{rm}{49}\RG{102.656}{-16.500}{7.9701}{rm}{50}\RG{125.795}{-16.500}{7.9701}{rm}{49}\RG{135.601}{-16.500}{7.9701}{rm}{50}\RR{48.686}{-4.653}{98.617}{0.436}\RR{212.786}{-4.653}{85.478}{0.436}%
}
\]
证明：\allowbreak{}\hskip.1em \RB{38.502}{8.685}{3.224}{%
\RG{0.000}{0.003}{10.9091}{rm}{40}\RG{4.243}{0.003}{10.9091}{mi}{86}\RG{11.203}{0.003}{10.9091}{mi}{59}\RG{16.061}{0.003}{10.9091}{sy}{8}\RG{24.546}{0.003}{10.9091}{mi}{59}\RG{29.394}{0.003}{10.9091}{sy}{14}\RG{34.849}{0.003}{10.9091}{rm}{41}%
}\hskip.1em\allowbreak{} 是实线性空间。

证明：用分式变换把这些运算化为实数的通常运算。定义

\[
\RB{208.434}{15.149}{10.707}{%
\RG{42.400}{7.383}{10.9091}{rm}{49}\RG{50.279}{7.383}{10.9091}{sy}{0}\RG{61.188}{7.383}{10.9091}{mi}{122}\RG{133.612}{7.383}{10.9091}{rm}{49}\RG{141.491}{7.383}{10.9091}{sy}{0}\RG{152.400}{7.383}{10.9091}{mi}{105}\RG{156.153}{7.383}{10.9091}{mi}{116}\RG{0.000}{0.003}{10.9091}{mi}{30}\RG{6.500}{0.003}{10.9091}{rm}{40}\RG{10.742}{0.003}{10.9091}{mi}{122}\RG{16.295}{0.003}{10.9091}{rm}{41}\RG{23.550}{0.003}{10.9091}{rm}{61}\RG{74.058}{0.003}{10.9091}{mi}{59}\RG{100.710}{0.003}{10.9091}{mi}{99}\RG{105.447}{0.003}{10.9091}{rm}{40}\RG{109.689}{0.003}{10.9091}{mi}{116}\RG{113.629}{0.003}{10.9091}{rm}{41}\RG{120.883}{0.003}{10.9091}{rm}{61}\RG{172.202}{0.003}{10.9091}{rm}{40}\RG{176.445}{0.003}{10.9091}{mi}{116}\RG{183.414}{0.003}{10.9091}{sy}{50}\RG{193.717}{0.003}{10.9091}{bb}{82}\RG{201.597}{0.003}{10.9091}{rm}{41}\RG{205.840}{0.003}{10.9091}{mi}{58}\RG{36.279}{-7.480}{10.9091}{mi}{105}\RG{40.037}{-7.480}{10.9091}{rm}{40}\RG{44.270}{-7.480}{10.9091}{rm}{49}\RG{52.146}{-7.480}{10.9091}{rm}{43}\RG{63.067}{-7.480}{10.9091}{mi}{122}\RG{68.620}{-7.480}{10.9091}{rm}{41}\RG{133.612}{-7.480}{10.9091}{rm}{49}\RG{141.488}{-7.480}{10.9091}{rm}{43}\RG{152.400}{-7.480}{10.9091}{mi}{105}\RG{156.153}{-7.480}{10.9091}{mi}{116}\RR{36.279}{2.730}{36.584}{0.436}\RR{133.612}{2.730}{26.486}{0.436}%
}
\]
若\allowbreak{}\hskip.1em \RB{27.775}{7.955}{0.932}{%
\RG{0.000}{0.004}{10.9091}{mi}{122}\RG{8.583}{0.004}{10.9091}{sy}{50}\RG{18.886}{0.004}{10.9091}{mi}{86}%
}\hskip.1em\allowbreak{} ，则\allowbreak{}\hskip.1em \RB{36.602}{8.686}{3.223}{%
\RG{0.000}{0.004}{10.9091}{mi}{122}\RG{0.655}{0.004}{10.9091}{rm}{22}\RG{8.583}{0.004}{10.9091}{rm}{61}\RG{20.081}{0.004}{10.9091}{rm}{49}\RG{25.536}{0.004}{10.9091}{rm}{47}\RG{31.007}{0.004}{10.9091}{mi}{122}%
}\hskip.1em\allowbreak{} ，从而

\[
\RB{130.081}{8.681}{18.091}{%
\RG{0.000}{-0.001}{10.9091}{rm}{49}\RG{7.879}{-0.001}{10.9091}{sy}{0}\RG{18.788}{-0.001}{10.9091}{mi}{122}\RG{41.277}{-0.001}{10.9091}{rm}{49}\RG{49.156}{-0.001}{10.9091}{sy}{0}\RG{60.065}{-0.001}{10.9091}{rm}{49}\RG{65.520}{-0.001}{10.9091}{rm}{47}\RG{70.975}{-0.001}{10.9091}{mi}{122}\RG{101.949}{-0.001}{10.9091}{rm}{49}\RG{109.828}{-0.001}{10.9091}{sy}{0}\RG{120.737}{-0.001}{10.9091}{mi}{122}\RG{28.567}{-7.381}{10.9091}{rm}{61}\RG{80.754}{-7.381}{10.9091}{rm}{61}\RG{92.269}{-7.381}{10.9091}{sy}{0}\RG{127.486}{-7.381}{10.9091}{mi}{58}\RG{0.000}{-14.864}{10.9091}{rm}{49}\RG{7.876}{-14.864}{10.9091}{rm}{43}\RG{18.788}{-14.864}{10.9091}{mi}{122}\RG{41.277}{-14.864}{10.9091}{rm}{49}\RG{49.154}{-14.864}{10.9091}{rm}{43}\RG{60.052}{-14.864}{10.9091}{rm}{49}\RG{65.507}{-14.864}{10.9091}{rm}{47}\RG{70.975}{-14.864}{10.9091}{mi}{122}\RG{101.949}{-14.864}{10.9091}{rm}{49}\RG{109.825}{-14.864}{10.9091}{rm}{43}\RG{120.737}{-14.864}{10.9091}{mi}{122}\RR{-1.196}{8.556}{26.732}{0.436}\RR{0.000}{-4.654}{24.341}{0.436}\RR{41.277}{-4.654}{35.250}{0.436}\RR{101.949}{-4.654}{24.341}{0.436}%
}
\]
故\allowbreak{}\hskip.1em \RB{70.516}{8.682}{3.227}{%
\RG{0.000}{0.000}{10.9091}{rm}{40}\RG{4.233}{0.000}{10.9091}{rm}{49}\RG{12.121}{0.000}{10.9091}{sy}{0}\RG{23.030}{0.000}{10.9091}{mi}{122}\RG{28.583}{0.000}{10.9091}{rm}{41}\RG{32.816}{0.000}{10.9091}{rm}{47}\RG{38.271}{0.000}{10.9091}{rm}{40}\RG{42.503}{0.000}{10.9091}{rm}{49}\RG{50.380}{0.000}{10.9091}{rm}{43}\RG{61.311}{0.000}{10.9091}{mi}{122}\RG{66.864}{0.000}{10.9091}{rm}{41}%
}\hskip.1em\allowbreak{} 为纯虚数，\allowbreak{}\hskip.1em \RB{42.040}{8.682}{3.227}{%
\RG{0.000}{0.000}{10.9091}{mi}{30}\RG{6.500}{0.000}{10.9091}{rm}{40}\RG{10.742}{0.000}{10.9091}{mi}{122}\RG{16.295}{0.000}{10.9091}{rm}{41}\RG{23.568}{0.000}{10.9091}{sy}{50}\RG{33.871}{0.000}{10.9091}{bb}{82}%
}\hskip.1em\allowbreak{} 。反之，对实数\allowbreak{}\hskip.1em \RB{4.100}{7.329}{0.620}{%
\RG{0.000}{0.000}{10.9091}{mi}{116}%
}\hskip.1em\allowbreak{} ，分母\allowbreak{}\hskip.1em \RB{46.532}{8.311}{2.856}{%
\RG{0.000}{0.000}{10.9091}{rm}{49}\RG{7.876}{0.000}{10.9091}{rm}{43}\RG{18.788}{0.000}{10.9091}{mi}{105}\RG{22.541}{0.000}{10.9091}{mi}{116}\RG{29.515}{0.000}{10.9091}{sy}{54}\RG{29.515}{0.000}{10.9091}{rm}{61}\RG{41.014}{0.000}{10.9091}{rm}{48}%
}\hskip.1em\allowbreak{} ，且

\[
\RB{171.823}{15.148}{8.887}{%
\RG{132.063}{7.382}{10.9091}{rm}{50}\RG{0.000}{0.002}{10.9091}{sy}{106}\RG{3.030}{0.002}{10.9091}{mi}{99}\RG{7.751}{0.002}{10.9091}{rm}{40}\RG{11.994}{0.002}{10.9091}{mi}{116}\RG{15.933}{0.002}{10.9091}{rm}{41}\RG{20.176}{0.002}{10.9091}{sy}{106}\RG{26.236}{0.002}{10.9091}{rm}{61}\RG{37.734}{0.002}{10.9091}{rm}{49}\RG{43.206}{0.002}{10.9091}{mi}{59}\RG{69.873}{0.002}{10.9091}{rm}{49}\RG{77.749}{0.002}{10.9091}{rm}{43}\RG{88.661}{0.002}{10.9091}{mi}{99}\RG{93.382}{0.002}{10.9091}{rm}{40}\RG{97.625}{0.002}{10.9091}{mi}{116}\RG{101.564}{0.002}{10.9091}{rm}{41}\RG{108.818}{0.002}{10.9091}{rm}{61}\RG{152.259}{0.002}{10.9091}{sy}{54}\RG{152.259}{0.002}{10.9091}{rm}{61}\RG{163.758}{0.002}{10.9091}{rm}{48}\RG{169.228}{0.002}{10.9091}{mi}{58}\RG{121.547}{-7.481}{10.9091}{rm}{49}\RG{129.423}{-7.481}{10.9091}{rm}{43}\RG{140.335}{-7.481}{10.9091}{mi}{105}\RG{144.088}{-7.481}{10.9091}{mi}{116}\RR{121.547}{2.730}{26.486}{0.436}%
}
\]
因此\allowbreak{}\hskip.1em \RB{39.367}{8.683}{3.226}{%
\RG{0.000}{0.001}{10.9091}{mi}{99}\RG{4.721}{0.001}{10.9091}{rm}{40}\RG{8.963}{0.001}{10.9091}{mi}{116}\RG{12.903}{0.001}{10.9091}{rm}{41}\RG{20.175}{0.001}{10.9091}{sy}{50}\RG{30.478}{0.001}{10.9091}{mi}{86}%
}\hskip.1em\allowbreak{} 。代入可验证\allowbreak{}\hskip.1em \RB{50.776}{8.683}{3.226}{%
\RG{0.000}{0.001}{10.9091}{mi}{30}\RG{6.500}{0.001}{10.9091}{rm}{40}\RG{10.743}{0.001}{10.9091}{mi}{99}\RG{15.464}{0.001}{10.9091}{rm}{40}\RG{19.706}{0.001}{10.9091}{mi}{116}\RG{23.646}{0.001}{10.9091}{rm}{41}\RG{27.879}{0.001}{10.9091}{rm}{41}\RG{35.134}{0.001}{10.9091}{rm}{61}\RG{46.676}{0.001}{10.9091}{mi}{116}%
}\hskip.1em\allowbreak{} 、\allowbreak{}\hskip.1em \RB{53.885}{8.683}{3.226}{%
\RG{0.000}{0.001}{10.9091}{mi}{99}\RG{4.721}{0.001}{10.9091}{rm}{40}\RG{8.964}{0.001}{10.9091}{mi}{30}\RG{15.464}{0.001}{10.9091}{rm}{40}\RG{19.706}{0.001}{10.9091}{mi}{122}\RG{25.259}{0.001}{10.9091}{rm}{41}\RG{29.492}{0.001}{10.9091}{rm}{41}\RG{36.747}{0.001}{10.9091}{rm}{61}\RG{48.290}{0.001}{10.9091}{mi}{122}%
}\hskip.1em\allowbreak{} ，所以\allowbreak{}\hskip.1em \RB{49.517}{8.072}{2.735}{%
\RG{0.000}{0.001}{10.9091}{mi}{30}\RG{9.530}{0.001}{10.9091}{rm}{58}\RG{15.590}{0.001}{10.9091}{mi}{86}\RG{27.409}{0.001}{10.9091}{sy}{33}\RG{41.348}{0.001}{10.9091}{bb}{82}%
}\hskip.1em\allowbreak{} 是双射，逆映射为\allowbreak{}\hskip.1em \RB{5.191}{5.323}{0.619}{%
\RG{0.000}{0.001}{10.9091}{mi}{99}%
}\hskip.1em\allowbreak{} 。

令\allowbreak{}\hskip.1em \RB{44.939}{8.678}{3.766}{%
\RG{0.000}{-0.004}{10.9091}{mi}{122}\RG{5.073}{-1.640}{7.9701}{mi}{106}\RG{12.486}{-0.004}{10.9091}{rm}{61}\RG{24.001}{-0.004}{10.9091}{mi}{99}\RG{28.722}{-0.004}{10.9091}{rm}{40}\RG{32.964}{-0.004}{10.9091}{mi}{116}\RG{36.903}{-1.640}{7.9701}{mi}{106}\RG{41.286}{-0.004}{10.9091}{rm}{41}%
}\hskip.1em\allowbreak{} 、\allowbreak{}\hskip.1em \RB{36.654}{8.678}{3.231}{%
\RG{0.000}{-0.004}{10.9091}{mi}{122}\RG{8.583}{-0.004}{10.9091}{rm}{61}\RG{20.099}{-0.004}{10.9091}{mi}{99}\RG{24.820}{-0.004}{10.9091}{rm}{40}\RG{29.062}{-0.004}{10.9091}{mi}{116}\RG{33.001}{-0.004}{10.9091}{rm}{41}%
}\hskip.1em\allowbreak{} ，其中\allowbreak{}\hskip.1em \RB{52.483}{7.969}{2.609}{%
\RG{0.000}{-0.004}{10.9091}{mi}{116}\RG{3.939}{-1.640}{7.9701}{rm}{49}\RG{8.672}{-0.004}{10.9091}{mi}{59}\RG{13.505}{-0.004}{10.9091}{mi}{116}\RG{17.459}{-1.640}{7.9701}{rm}{50}\RG{22.192}{-0.004}{10.9091}{mi}{59}\RG{27.025}{-0.004}{10.9091}{mi}{116}\RG{34.010}{-0.004}{10.9091}{sy}{50}\RG{44.313}{-0.004}{10.9091}{bb}{82}%
}\hskip.1em\allowbreak{} 。原题两个分母分别为

\[
\RB{184.535}{47.600}{8.892}{%
\RG{107.795}{46.663}{10.9092}{ex}{0}\RG{174.503}{46.663}{10.9092}{ex}{1}\RG{102.341}{37.827}{10.9091}{rm}{52}\RG{112.795}{37.827}{10.9091}{rm}{49}\RG{120.671}{37.827}{10.9091}{rm}{43}\RG{131.583}{37.827}{10.9091}{mi}{105}\RG{135.342}{37.827}{10.9091}{rm}{40}\RG{139.584}{37.827}{10.9091}{mi}{116}\RG{150.680}{37.827}{10.9091}{rm}{43}\RG{161.589}{37.827}{10.9091}{mi}{116}\RG{170.261}{37.827}{10.9091}{rm}{41}\RG{143.523}{36.191}{7.9701}{rm}{49}\RG{165.528}{36.191}{7.9701}{rm}{50}\RG{0.801}{29.754}{10.9091}{rm}{51}\RG{8.677}{29.754}{10.9091}{rm}{43}\RG{19.589}{29.754}{10.9091}{mi}{122}\RG{31.818}{29.754}{10.9091}{rm}{43}\RG{42.727}{29.754}{10.9091}{mi}{122}\RG{54.957}{29.754}{10.9091}{sy}{0}\RG{65.867}{29.754}{10.9091}{mi}{122}\RG{75.673}{29.754}{10.9091}{mi}{122}\RG{88.509}{29.754}{10.9091}{rm}{61}\RG{181.821}{29.754}{10.9091}{mi}{59}\RG{24.662}{28.118}{7.9701}{rm}{49}\RG{47.801}{28.118}{7.9701}{rm}{50}\RG{70.941}{28.118}{7.9701}{rm}{49}\RG{80.746}{28.118}{7.9701}{rm}{50}\RG{101.219}{22.271}{10.9091}{rm}{40}\RG{105.452}{22.271}{10.9091}{rm}{49}\RG{113.328}{22.271}{10.9091}{rm}{43}\RG{124.249}{22.271}{10.9091}{mi}{105}\RG{128.002}{22.271}{10.9091}{mi}{116}\RG{136.679}{22.271}{10.9091}{rm}{41}\RG{140.912}{22.271}{10.9091}{rm}{40}\RG{145.145}{22.271}{10.9091}{rm}{49}\RG{153.021}{22.271}{10.9091}{rm}{43}\RG{163.952}{22.271}{10.9091}{mi}{105}\RG{167.705}{22.271}{10.9091}{mi}{116}\RG{176.382}{22.271}{10.9091}{rm}{41}\RG{131.947}{20.635}{7.9701}{rm}{49}\RG{171.650}{20.635}{7.9701}{rm}{50}\RG{101.219}{7.377}{10.9091}{rm}{50}\RG{106.673}{7.377}{10.9091}{rm}{40}\RG{110.906}{7.377}{10.9091}{rm}{49}\RG{118.783}{7.377}{10.9091}{rm}{43}\RG{129.704}{7.377}{10.9091}{mi}{105}\RG{133.457}{7.377}{10.9091}{mi}{107}\RG{139.468}{7.377}{10.9091}{mi}{116}\RG{143.424}{7.377}{10.9091}{rm}{41}\RG{0.000}{-0.003}{10.9091}{rm}{40}\RG{4.233}{-0.003}{10.9091}{rm}{49}\RG{12.109}{-0.003}{10.9091}{rm}{43}\RG{23.030}{-0.003}{10.9091}{mi}{107}\RG{29.053}{-0.003}{10.9091}{rm}{41}\RG{35.708}{-0.003}{10.9091}{rm}{43}\RG{46.606}{-0.003}{10.9091}{rm}{40}\RG{50.839}{-0.003}{10.9091}{rm}{49}\RG{58.750}{-0.003}{10.9091}{sy}{0}\RG{69.660}{-0.003}{10.9091}{mi}{107}\RG{75.683}{-0.003}{10.9091}{rm}{41}\RG{79.925}{-0.003}{10.9091}{mi}{122}\RG{88.509}{-0.003}{10.9091}{rm}{61}\RG{148.862}{-0.003}{10.9091}{mi}{58}\RG{111.200}{-7.486}{10.9091}{rm}{49}\RG{119.076}{-7.486}{10.9091}{rm}{43}\RG{129.988}{-7.486}{10.9091}{mi}{105}\RG{133.741}{-7.486}{10.9091}{mi}{116}\RR{101.219}{32.481}{79.407}{0.436}\RR{101.219}{2.724}{46.448}{0.436}%
}
\]
它们都不为零。化简分子和分母，得到

\[
\RB{180.227}{8.678}{3.231}{%
\RG{0.000}{-0.004}{10.9091}{mi}{122}\RG{12.230}{-0.004}{10.9091}{sy}{8}\RG{23.139}{-0.004}{10.9091}{mi}{122}\RG{35.975}{-0.004}{10.9091}{rm}{61}\RG{47.490}{-0.004}{10.9091}{mi}{99}\RG{52.211}{-0.004}{10.9091}{rm}{40}\RG{56.453}{-0.004}{10.9091}{mi}{116}\RG{67.550}{-0.004}{10.9091}{rm}{43}\RG{78.459}{-0.004}{10.9091}{mi}{116}\RG{87.131}{-0.004}{10.9091}{rm}{41}\RG{91.373}{-0.004}{10.9091}{mi}{59}\RG{118.024}{-0.004}{10.9091}{mi}{107}\RG{126.487}{-0.004}{10.9091}{sy}{14}\RG{134.366}{-0.004}{10.9091}{mi}{122}\RG{142.949}{-0.004}{10.9091}{rm}{61}\RG{154.464}{-0.004}{10.9091}{mi}{99}\RG{159.185}{-0.004}{10.9091}{rm}{40}\RG{163.428}{-0.004}{10.9091}{mi}{107}\RG{169.439}{-0.004}{10.9091}{mi}{116}\RG{173.390}{-0.004}{10.9091}{rm}{41}\RG{177.632}{-0.004}{10.9091}{mi}{58}\RG{5.073}{-1.640}{7.9701}{rm}{49}\RG{28.212}{-1.640}{7.9701}{rm}{50}\RG{60.393}{-1.640}{7.9701}{rm}{49}\RG{82.399}{-1.640}{7.9701}{rm}{50}%
}
\]
这既证明两种运算在\allowbreak{}\hskip.1em \RB{8.889}{7.948}{0.743}{%
\RG{0.000}{-0.003}{10.9091}{mi}{86}%
}\hskip.1em\allowbreak{} 中封闭，又给出

\[
\RB{232.622}{8.681}{3.228}{%
\RG{0.000}{-0.001}{10.9091}{mi}{30}\RG{6.500}{-0.001}{10.9091}{rm}{40}\RG{10.743}{-0.001}{10.9091}{mi}{122}\RG{22.972}{-0.001}{10.9091}{sy}{8}\RG{33.881}{-0.001}{10.9091}{mi}{122}\RG{43.687}{-0.001}{10.9091}{rm}{41}\RG{50.942}{-0.001}{10.9091}{rm}{61}\RG{62.476}{-0.001}{10.9091}{mi}{30}\RG{68.976}{-0.001}{10.9091}{rm}{40}\RG{73.218}{-0.001}{10.9091}{mi}{122}\RG{83.024}{-0.001}{10.9091}{rm}{41}\RG{89.678}{-0.001}{10.9091}{rm}{43}\RG{100.600}{-0.001}{10.9091}{mi}{30}\RG{107.100}{-0.001}{10.9091}{rm}{40}\RG{111.342}{-0.001}{10.9091}{mi}{122}\RG{121.148}{-0.001}{10.9091}{rm}{41}\RG{125.390}{-0.001}{10.9091}{mi}{59}\RG{152.041}{-0.001}{10.9091}{mi}{30}\RG{158.557}{-0.001}{10.9091}{rm}{40}\RG{162.799}{-0.001}{10.9091}{mi}{107}\RG{171.246}{-0.001}{10.9091}{sy}{14}\RG{179.126}{-0.001}{10.9091}{mi}{122}\RG{184.679}{-0.001}{10.9091}{rm}{41}\RG{191.933}{-0.001}{10.9091}{rm}{61}\RG{203.467}{-0.001}{10.9091}{mi}{107}\RG{209.478}{-0.001}{10.9091}{mi}{30}\RG{215.989}{-0.001}{10.9091}{rm}{40}\RG{220.232}{-0.001}{10.9091}{mi}{122}\RG{225.785}{-0.001}{10.9091}{rm}{41}\RG{230.027}{-0.001}{10.9091}{mi}{58}\RG{15.816}{-1.638}{7.9701}{rm}{49}\RG{38.955}{-1.638}{7.9701}{rm}{50}\RG{78.292}{-1.638}{7.9701}{rm}{49}\RG{116.415}{-1.638}{7.9701}{rm}{50}%
}
\]
由\allowbreak{}\hskip.1em \RB{6.751}{8.071}{2.736}{%
\RG{0.000}{0.000}{10.9091}{mi}{30}%
}\hskip.1em\allowbreak{} 的单射性，实数加法的交换律、结合律立即转为\allowbreak{}\hskip.1em \RB{8.376}{6.860}{1.405}{%
\RG{0.000}{0.000}{10.9091}{sy}{8}%
}\hskip.1em\allowbreak{} 的交换律、结合律。零向量和加法逆元分别是

\[
\RB{193.445}{8.682}{3.227}{%
\RG{0.000}{-0.000}{10.9091}{rm}{48}\RG{15.796}{-0.000}{10.9091}{rm}{61}\RG{27.311}{-0.000}{10.9091}{mi}{99}\RG{32.032}{-0.000}{10.9091}{rm}{40}\RG{36.265}{-0.000}{10.9091}{rm}{48}\RG{41.720}{-0.000}{10.9091}{rm}{41}\RG{48.974}{-0.000}{10.9091}{rm}{61}\RG{60.483}{-0.000}{10.9091}{rm}{49}\RG{65.972}{-0.000}{10.9091}{mi}{59}\RG{92.639}{-0.000}{10.9091}{sy}{0}\RG{108.435}{-0.000}{10.9091}{mi}{122}\RG{117.019}{-0.000}{10.9091}{rm}{61}\RG{128.534}{-0.000}{10.9091}{mi}{99}\RG{133.255}{-0.000}{10.9091}{rm}{40}\RG{137.497}{-0.000}{10.9091}{sy}{0}\RG{145.982}{-0.000}{10.9091}{mi}{30}\RG{152.482}{-0.000}{10.9091}{rm}{40}\RG{156.725}{-0.000}{10.9091}{mi}{122}\RG{162.278}{-0.000}{10.9091}{rm}{41}\RG{166.511}{-0.000}{10.9091}{rm}{41}\RG{173.766}{-0.000}{10.9091}{rm}{61}\RG{185.309}{-0.000}{10.9091}{mi}{122}\RG{185.929}{-0.000}{10.9091}{rm}{22}\RG{190.850}{-0.000}{10.9091}{mi}{58}\RG{5.454}{-1.689}{7.9701}{mi}{86}\RG{101.124}{-1.689}{7.9701}{mi}{86}%
}
\]
同样，对实数\allowbreak{}\hskip.1em \RB{15.620}{8.073}{2.603}{%
\RG{0.000}{0.002}{10.9091}{mi}{97}\RG{5.760}{0.002}{10.9091}{mi}{59}\RG{10.593}{0.002}{10.9091}{mi}{98}%
}\hskip.1em\allowbreak{} ，下列等式经\allowbreak{}\hskip.1em \RB{6.751}{8.073}{2.734}{%
\RG{0.000}{0.002}{10.9091}{mi}{30}%
}\hskip.1em\allowbreak{} 映射后分别化为实数中的对应恒等式，故均成立：

\[
\RB{351.387}{26.845}{3.228}{%
\RG{26.812}{18.163}{10.9091}{rm}{49}\RG{34.691}{18.163}{10.9091}{sy}{14}\RG{42.569}{18.163}{10.9091}{mi}{122}\RG{51.153}{18.163}{10.9091}{rm}{61}\RG{62.669}{18.163}{10.9091}{mi}{122}\RG{68.210}{18.163}{10.9091}{mi}{59}\RG{207.101}{18.163}{10.9091}{mi}{97}\RG{215.292}{18.163}{10.9091}{sy}{14}\RG{223.170}{18.163}{10.9091}{rm}{40}\RG{227.413}{18.163}{10.9091}{mi}{98}\RG{234.519}{18.163}{10.9091}{sy}{14}\RG{242.398}{18.163}{10.9091}{mi}{122}\RG{247.951}{18.163}{10.9091}{rm}{41}\RG{255.206}{18.163}{10.9091}{rm}{61}\RG{266.715}{18.163}{10.9091}{rm}{40}\RG{270.981}{18.163}{10.9091}{mi}{97}\RG{276.741}{18.163}{10.9091}{mi}{98}\RG{281.429}{18.163}{10.9091}{rm}{41}\RG{288.096}{18.163}{10.9091}{sy}{14}\RG{295.975}{18.163}{10.9091}{mi}{122}\RG{301.517}{18.163}{10.9091}{mi}{59}\RG{0.000}{-0.001}{10.9091}{rm}{40}\RG{4.242}{-0.001}{10.9091}{mi}{97}\RG{12.433}{-0.001}{10.9091}{rm}{43}\RG{23.342}{-0.001}{10.9091}{mi}{98}\RG{28.024}{-0.001}{10.9091}{rm}{41}\RG{34.691}{-0.001}{10.9091}{sy}{14}\RG{42.569}{-0.001}{10.9091}{mi}{122}\RG{51.153}{-0.001}{10.9091}{rm}{61}\RG{62.652}{-0.001}{10.9091}{rm}{40}\RG{66.911}{-0.001}{10.9091}{mi}{97}\RG{75.102}{-0.001}{10.9091}{sy}{14}\RG{82.981}{-0.001}{10.9091}{mi}{122}\RG{88.534}{-0.001}{10.9091}{rm}{41}\RG{95.200}{-0.001}{10.9091}{sy}{8}\RG{106.109}{-0.001}{10.9091}{rm}{40}\RG{110.352}{-0.001}{10.9091}{mi}{98}\RG{117.458}{-0.001}{10.9091}{sy}{14}\RG{125.337}{-0.001}{10.9091}{mi}{122}\RG{130.890}{-0.001}{10.9091}{rm}{41}\RG{135.132}{-0.001}{10.9091}{mi}{59}\RG{194.694}{-0.001}{10.9091}{mi}{97}\RG{202.885}{-0.001}{10.9091}{sy}{14}\RG{210.764}{-0.001}{10.9091}{rm}{40}\RG{215.006}{-0.001}{10.9091}{mi}{122}\RG{227.236}{-0.001}{10.9091}{sy}{8}\RG{238.145}{-0.001}{10.9091}{mi}{122}\RG{247.951}{-0.001}{10.9091}{rm}{41}\RG{255.206}{-0.001}{10.9091}{rm}{61}\RG{266.715}{-0.001}{10.9091}{rm}{40}\RG{270.981}{-0.001}{10.9091}{mi}{97}\RG{279.172}{-0.001}{10.9091}{sy}{14}\RG{287.050}{-0.001}{10.9091}{mi}{122}\RG{296.856}{-0.001}{10.9091}{rm}{41}\RG{303.524}{-0.001}{10.9091}{sy}{8}\RG{314.433}{-0.001}{10.9091}{rm}{40}\RG{318.675}{-0.001}{10.9091}{mi}{97}\RG{326.866}{-0.001}{10.9091}{sy}{14}\RG{334.745}{-0.001}{10.9091}{mi}{122}\RG{344.550}{-0.001}{10.9091}{rm}{41}\RG{348.793}{-0.001}{10.9091}{mi}{58}\RG{220.080}{-1.637}{7.9701}{rm}{49}\RG{243.218}{-1.637}{7.9701}{rm}{50}\RG{292.124}{-1.637}{7.9701}{rm}{49}\RG{339.818}{-1.637}{7.9701}{rm}{50}%
}
\]
线性空间的各条公理均已满足，所求结论成立。\allowbreak{}\hskip.1em \RB{8.366}{7.995}{0.500}{%
\RG{0.000}{0.000}{10.9091}{am}{3}%
}\hskip.1em\allowbreak{}

方法小结：这些分式运算来自实直线上的加法和数乘。双射\allowbreak{}\hskip.1em \RB{6.751}{8.068}{2.739}{%
\RG{0.000}{-0.003}{10.9091}{mi}{30}%
}\hskip.1em\allowbreak{} 同时负责验证分母非零、运算封闭和各条公理；该空间实际上为一维，且\allowbreak{}\hskip.1em \RB{22.856}{8.684}{3.225}{%
\RG{0.000}{0.002}{10.9091}{sy}{102}\RG{5.455}{0.002}{10.9091}{sy}{0}\RG{13.940}{0.002}{10.9091}{mi}{105}\RG{17.698}{0.002}{10.9091}{sy}{103}%
}\hskip.1em\allowbreak{} 是它的一组基。

题目来源：谢启鸿《复旦大学高等代数习题集（2023版）》，第124页。

% Original page 5.
\section*{2022A09}
设\allowbreak{}\hskip.1em \RB{8.366}{8.310}{0.501}{%
\RG{0.000}{-0.001}{10.9091}{mi}{65}%
}\hskip.1em\allowbreak{} 为\allowbreak{}\hskip.1em \RB{6.729}{5.321}{0.621}{%
\RG{0.000}{-0.001}{10.9091}{mi}{110}%
}\hskip.1em\allowbreak{} 阶方阵，其主对角元素全为\allowbreak{}\hskip.1em \RB{5.518}{7.765}{0.741}{%
\RG{0.000}{-0.001}{10.9091}{rm}{48}%
}\hskip.1em\allowbreak{} ，其余元素为\allowbreak{}\hskip.1em \RB{5.071}{7.765}{0.501}{%
\RG{0.000}{-0.001}{10.9091}{rm}{49}%
}\hskip.1em\allowbreak{} 或\allowbreak{}\hskip.1em \RB{13.556}{7.765}{0.501}{%
\RG{0.000}{-0.001}{10.9091}{sy}{0}\RG{8.485}{-0.001}{10.9091}{rm}{49}%
}\hskip.1em\allowbreak{} 。证明：

\[
\RB{80.699}{9.293}{3.227}{%
\RT{0.000}{-0.000}{10.9091}{r}\RT{4.276}{-0.000}{10.9091}{a}\RT{9.731}{-0.000}{10.9091}{n}\RT{15.796}{-0.000}{10.9091}{k}\RG{21.556}{-0.000}{10.9091}{rm}{40}\RG{25.798}{-0.000}{10.9091}{mi}{65}\RG{33.980}{-0.000}{10.9091}{rm}{41}\RG{41.253}{-0.000}{10.9091}{sy}{21}\RG{52.768}{-0.000}{10.9091}{mi}{110}\RG{61.741}{-0.000}{10.9091}{sy}{0}\RG{72.650}{-0.000}{10.9091}{rm}{49}\RG{78.105}{-0.000}{10.9091}{mi}{58}%
}
\]
证明：先考虑\allowbreak{}\hskip.1em \RB{6.729}{5.323}{0.619}{%
\RG{0.000}{0.001}{10.9091}{mi}{110}%
}\hskip.1em\allowbreak{} 为偶数的情形。由于矩阵元素都是整数，可以将它们模\allowbreak{}\hskip.1em \RB{5.398}{7.767}{0.500}{%
\RG{0.000}{0.001}{10.9091}{rm}{50}%
}\hskip.1em\allowbreak{} 化简，得到二元域\allowbreak{}\hskip.1em \RB{10.969}{7.974}{2.135}{%
\RG{0.000}{0.001}{10.9091}{bb}{70}\RG{6.667}{-1.635}{7.9701}{rm}{50}%
}\hskip.1em\allowbreak{} 上的矩阵\allowbreak{}\hskip.1em \RB{8.366}{9.689}{0.504}{%
\RG{0.000}{-0.004}{10.9091}{mi}{65}\RG{2.879}{2.753}{10.9091}{rm}{22}%
}\hskip.1em\allowbreak{} 。因为

\[
\RB{85.746}{9.294}{3.226}{%
\RG{0.000}{0.001}{10.9091}{rm}{49}\RG{8.485}{0.001}{10.9091}{sy}{17}\RG{19.983}{0.001}{10.9091}{sy}{0}\RG{28.485}{0.001}{10.9091}{rm}{49}\RG{44.838}{0.001}{10.9091}{rm}{40}\RT{49.091}{0.001}{10.9091}{m}\RT{58.179}{0.001}{10.9091}{o}\RT{63.633}{0.001}{10.9091}{d}\RG{73.335}{0.001}{10.9091}{rm}{50}\RG{78.789}{0.001}{10.9091}{rm}{41}\RG{83.032}{0.001}{10.9091}{mi}{59}%
}
\]
所以\allowbreak{}\hskip.1em \RB{8.366}{9.697}{0.500}{%
\RG{0.000}{0.003}{10.9091}{mi}{65}\RG{2.879}{2.761}{10.9091}{rm}{22}%
}\hskip.1em\allowbreak{} 的对角元素全为\allowbreak{}\hskip.1em \RB{5.518}{7.769}{0.737}{%
\RG{0.000}{0.003}{10.9091}{rm}{48}%
}\hskip.1em\allowbreak{} ，其余元素全为\allowbreak{}\hskip.1em \RB{5.071}{7.769}{0.500}{%
\RG{0.000}{0.003}{10.9091}{rm}{49}%
}\hskip.1em\allowbreak{} 。

设\allowbreak{}\hskip.1em \RB{75.756}{8.680}{3.229}{%
\RG{0.000}{-0.002}{10.9091}{mi}{120}\RG{9.265}{-0.002}{10.9091}{rm}{61}\RG{20.763}{-0.002}{10.9091}{rm}{40}\RG{25.023}{-0.002}{10.9091}{mi}{120}\RG{31.257}{-1.638}{7.9701}{rm}{49}\RG{35.990}{-0.002}{10.9091}{mi}{59}\RG{40.823}{-0.002}{10.9091}{mi}{58}\RG{45.667}{-0.002}{10.9091}{mi}{58}\RG{50.500}{-0.002}{10.9091}{mi}{58}\RG{55.343}{-0.002}{10.9091}{mi}{59}\RG{60.187}{-0.002}{10.9091}{mi}{120}\RG{66.468}{-1.638}{7.9701}{mi}{110}\RG{72.104}{-0.002}{10.9091}{rm}{41}%
}\hskip.1em\allowbreak{} T\allowbreak{}\hskip.1em \RB{22.348}{7.972}{3.425}{%
\RG{0.000}{-0.002}{10.9091}{sy}{50}\RG{10.303}{-0.002}{10.9091}{bb}{70}\RG{16.970}{3.957}{7.9701}{mi}{110}\RG{16.970}{-2.925}{7.9701}{rm}{50}%
}\hskip.1em\allowbreak{} 满足\allowbreak{}\hskip.1em \RB{34.464}{9.692}{0.742}{%
\RG{0.000}{-0.002}{10.9091}{mi}{65}\RG{2.879}{2.756}{10.9091}{rm}{22}\RG{8.182}{-0.002}{10.9091}{mi}{120}\RG{17.447}{-0.002}{10.9091}{rm}{61}\RG{28.945}{-0.002}{10.9091}{rm}{48}%
}\hskip.1em\allowbreak{} ，并记\allowbreak{}\hskip.1em \RB{81.632}{6.858}{2.218}{%
\RG{0.000}{-0.002}{10.9091}{mi}{115}\RG{8.143}{-0.002}{10.9091}{rm}{61}\RG{19.659}{-0.002}{10.9091}{mi}{120}\RG{25.893}{-1.638}{7.9701}{rm}{49}\RG{33.050}{-0.002}{10.9091}{rm}{43}\RG{43.959}{-0.002}{10.9091}{sy}{1}\RG{48.792}{-0.002}{10.9091}{sy}{1}\RG{53.636}{-0.002}{10.9091}{sy}{1}\RG{59.111}{-0.002}{10.9091}{rm}{43}\RG{70.021}{-0.002}{10.9091}{mi}{120}\RG{76.255}{-1.638}{7.9701}{mi}{110}%
}\hskip.1em\allowbreak{} 。第\allowbreak{}\hskip.1em \RB{3.696}{7.709}{0.622}{%
\RG{0.000}{-0.002}{10.9091}{mi}{105}%
}\hskip.1em\allowbreak{} 个方程给出

\[
\RB{50.908}{10.863}{15.568}{%
\RG{0.000}{10.363}{10.9092}{ex}{88}\RG{17.576}{-0.001}{10.9091}{mi}{120}\RG{31.223}{-0.001}{10.9091}{rm}{61}\RG{42.722}{-0.001}{10.9091}{rm}{48}\RG{48.193}{-0.001}{10.9091}{mi}{59}\RG{23.811}{-1.637}{7.9701}{mi}{106}\RG{1.202}{-13.354}{7.9701}{mi}{106}\RG{5.086}{-13.354}{7.9701}{sy}{54}\RG{5.086}{-13.354}{7.9701}{rm}{61}\RG{11.673}{-13.354}{7.9701}{mi}{105}%
}
\]
即\allowbreak{}\hskip.1em \RB{48.110}{7.770}{2.212}{%
\RG{0.000}{0.004}{10.9091}{mi}{115}\RG{7.538}{0.004}{10.9091}{sy}{0}\RG{18.447}{0.004}{10.9091}{mi}{120}\RG{24.682}{-1.632}{7.9701}{mi}{105}\RG{31.093}{0.004}{10.9091}{rm}{61}\RG{42.592}{0.004}{10.9091}{rm}{48}%
}\hskip.1em\allowbreak{} 。所以\allowbreak{}\hskip.1em \RB{29.243}{5.326}{2.212}{%
\RG{0.000}{0.004}{10.9091}{mi}{120}\RG{6.235}{-1.632}{7.9701}{mi}{105}\RG{12.646}{0.004}{10.9091}{rm}{61}\RG{24.161}{0.004}{10.9091}{mi}{115}%
}\hskip.1em\allowbreak{} 对所有\allowbreak{}\hskip.1em \RB{3.696}{7.715}{0.616}{%
\RG{0.000}{0.004}{10.9091}{mi}{105}%
}\hskip.1em\allowbreak{} 成立，再将这些等式相加得

\allowbreak{}\hskip.1em \RB{51.368}{7.771}{0.735}{%
\RG{0.000}{0.005}{10.9091}{mi}{115}\RG{8.144}{0.005}{10.9091}{rm}{61}\RG{19.659}{0.005}{10.9091}{mi}{110}\RG{26.205}{0.005}{10.9091}{mi}{115}\RG{34.351}{0.005}{10.9091}{rm}{61}\RG{45.850}{0.005}{10.9091}{rm}{48}%
}\hskip.1em\allowbreak{} （\allowbreak{}\hskip.1em \RB{6.729}{5.327}{0.615}{%
\RG{0.000}{0.005}{10.9091}{mi}{110}%
}\hskip.1em\allowbreak{} 为偶数，在\allowbreak{}\hskip.1em \RB{10.968}{7.978}{2.131}{%
\RG{0.000}{0.005}{10.9091}{bb}{70}\RG{6.666}{-1.631}{7.9701}{rm}{50}%
}\hskip.1em\allowbreak{} 中\allowbreak{}\hskip.1em \RB{26.596}{7.771}{0.735}{%
\RG{0.000}{0.005}{10.9091}{mi}{110}\RG{9.579}{0.005}{10.9091}{rm}{61}\RG{21.077}{0.005}{10.9091}{rm}{48}%
}\hskip.1em\allowbreak{} ）\allowbreak{}\hskip.1em \RB{2.595}{1.661}{0.500}{%
\RG{0.000}{0.005}{10.9091}{mi}{58}%
}\hskip.1em\allowbreak{}

于是\allowbreak{}\hskip.1em \RB{26.282}{7.762}{0.744}{%
\RG{0.000}{-0.004}{10.9091}{mi}{120}\RG{9.265}{-0.004}{10.9091}{rm}{61}\RG{20.763}{-0.004}{10.9091}{rm}{48}%
}\hskip.1em\allowbreak{} ，故\allowbreak{}\hskip.1em \RB{8.366}{9.690}{0.504}{%
\RG{0.000}{-0.004}{10.9091}{mi}{65}\RG{2.878}{2.754}{10.9091}{rm}{22}%
}\hskip.1em\allowbreak{} 可逆。因此

\[
\RB{136.539}{9.691}{3.230}{%
\RG{59.549}{2.755}{10.9091}{rm}{22}\RT{0.000}{-0.003}{10.9091}{d}\RT{6.066}{-0.003}{10.9091}{e}\RT{10.909}{-0.003}{10.9091}{t}\RG{16.971}{-0.003}{10.9091}{mi}{65}\RG{28.183}{-0.003}{10.9091}{sy}{17}\RT{39.699}{-0.003}{10.9091}{d}\RT{45.765}{-0.003}{10.9091}{e}\RT{50.609}{-0.003}{10.9091}{t}\RG{56.669}{-0.003}{10.9091}{mi}{65}\RG{67.883}{-0.003}{10.9091}{rm}{61}\RG{79.381}{-0.003}{10.9091}{rm}{49}\RG{95.745}{-0.003}{10.9091}{rm}{40}\RT{100.004}{-0.003}{10.9091}{m}\RT{109.091}{-0.003}{10.9091}{o}\RT{114.546}{-0.003}{10.9091}{d}\RG{124.247}{-0.003}{10.9091}{rm}{50}\RG{129.702}{-0.003}{10.9091}{rm}{41}\RG{133.944}{-0.003}{10.9091}{mi}{58}%
}
\]
这里行列式按模\allowbreak{}\hskip.1em \RB{5.398}{7.763}{0.502}{%
\RG{0.000}{-0.002}{10.9091}{rm}{50}%
}\hskip.1em\allowbreak{} 化简与先计算再化简相容，因为它是矩阵元素的整系数多项式。于是det\allowbreak{}\hskip.1em \RB{8.366}{8.309}{0.502}{%
\RG{0.000}{-0.002}{10.9091}{mi}{65}%
}\hskip.1em\allowbreak{} 为奇整数，特别地不为零，从而

rank\allowbreak{}\hskip.1em \RB{70.551}{8.680}{3.229}{%
\RG{0.000}{-0.002}{10.9091}{rm}{40}\RG{4.243}{-0.002}{10.9091}{mi}{65}\RG{12.425}{-0.002}{10.9091}{rm}{41}\RG{19.680}{-0.002}{10.9091}{rm}{61}\RG{31.213}{-0.002}{10.9091}{mi}{110}\RG{59.579}{-0.002}{10.9091}{rm}{40}\RG{63.822}{-0.002}{10.9091}{mi}{110}%
}\hskip.1em\allowbreak{} 为偶数\allowbreak{}\hskip.1em \RB{6.838}{8.680}{3.229}{%
\RG{0.000}{-0.002}{10.9091}{rm}{41}\RG{4.243}{-0.002}{10.9091}{mi}{58}%
}\hskip.1em\allowbreak{}

当\allowbreak{}\hskip.1em \RB{26.951}{7.767}{1.994}{%
\RG{0.000}{0.001}{10.9091}{mi}{110}\RG{9.764}{0.001}{10.9091}{sy}{21}\RG{21.465}{0.001}{10.9091}{rm}{51}%
}\hskip.1em\allowbreak{} 为奇数时，取\allowbreak{}\hskip.1em \RB{8.366}{8.312}{0.500}{%
\RG{0.000}{0.001}{10.9091}{mi}{65}%
}\hskip.1em\allowbreak{} 的前\allowbreak{}\hskip.1em \RB{25.101}{7.767}{0.619}{%
\RG{0.000}{0.001}{10.9091}{mi}{110}\RG{9.047}{0.001}{10.9091}{sy}{0}\RG{20.030}{0.001}{10.9091}{rm}{49}%
}\hskip.1em\allowbreak{} 行与前\allowbreak{}\hskip.1em \RB{25.101}{7.767}{0.619}{%
\RG{0.000}{0.001}{10.9091}{mi}{110}\RG{9.046}{0.001}{10.9091}{sy}{0}\RG{20.030}{0.001}{10.9091}{rm}{49}%
}\hskip.1em\allowbreak{} 列组成的主子矩阵\allowbreak{}\hskip.1em \RB{12.579}{8.312}{2.303}{%
\RG{0.000}{0.001}{10.9091}{mi}{65}\RG{8.182}{-1.636}{7.9701}{rm}{48}%
}\hskip.1em\allowbreak{} 。它的阶数为偶数，并满足同样的元素条件，故由上面的结论，

\[
\RB{52.480}{9.290}{2.859}{%
\RT{0.000}{-0.003}{10.9091}{d}\RT{6.066}{-0.003}{10.9091}{e}\RT{10.909}{-0.003}{10.9091}{t}\RG{16.971}{-0.003}{10.9091}{mi}{65}\RG{32.916}{-0.003}{10.9091}{sy}{54}\RG{32.916}{-0.003}{10.9091}{rm}{61}\RG{44.415}{-0.003}{10.9091}{rm}{48}\RG{49.885}{-0.003}{10.9091}{mi}{58}\RG{25.153}{-1.640}{7.9701}{rm}{48}%
}
\]
这给出一个非零的\allowbreak{}\hskip.1em \RB{24.823}{7.764}{0.622}{%
\RG{0.000}{-0.002}{10.9091}{mi}{110}\RG{8.908}{-0.002}{10.9091}{sy}{0}\RG{19.752}{-0.002}{10.9091}{rm}{49}%
}\hskip.1em\allowbreak{} 阶子式，因此rank\allowbreak{}\hskip.1em \RB{56.035}{8.680}{3.229}{%
\RG{0.000}{-0.002}{10.9091}{rm}{40}\RG{4.242}{-0.002}{10.9091}{mi}{65}\RG{12.424}{-0.002}{10.9091}{rm}{41}\RG{19.697}{-0.002}{10.9091}{sy}{21}\RG{31.212}{-0.002}{10.9091}{mi}{110}\RG{40.119}{-0.002}{10.9091}{sy}{0}\RG{50.964}{-0.002}{10.9091}{rm}{49}%
}\hskip.1em\allowbreak{} 。当\allowbreak{}\hskip.1em \RB{26.147}{7.764}{0.622}{%
\RG{0.000}{-0.002}{10.9091}{mi}{110}\RG{9.578}{-0.002}{10.9091}{rm}{61}\RG{21.076}{-0.002}{10.9091}{rm}{49}%
}\hskip.1em\allowbreak{} 时，\allowbreak{}\hskip.1em \RB{36.051}{8.680}{3.229}{%
\RG{0.000}{-0.002}{10.9091}{mi}{65}\RG{11.212}{-0.002}{10.9091}{rm}{61}\RG{22.710}{-0.002}{10.9091}{rm}{40}\RG{26.943}{-0.002}{10.9091}{rm}{48}\RG{32.398}{-0.002}{10.9091}{rm}{41}%
}\hskip.1em\allowbreak{} ，结论也成立。\allowbreak{}\hskip.1em \RB{8.366}{7.993}{0.502}{%
\RG{0.000}{-0.002}{10.9091}{am}{3}%
}\hskip.1em\allowbreak{}补充：奇数阶时此下界可以达到。取\allowbreak{}\hskip.1em \RB{32.630}{7.769}{3.760}{%
\RG{0.000}{0.003}{10.9091}{mi}{97}\RG{5.766}{-1.634}{7.9701}{mi}{105}\RG{8.643}{-1.634}{7.9701}{mi}{106}\RG{16.061}{0.003}{10.9091}{rm}{61}\RG{27.559}{0.003}{10.9091}{rm}{49}%
}\hskip.1em\allowbreak{} （\allowbreak{}\hskip.1em \RB{23.115}{7.714}{2.733}{%
\RG{0.000}{0.003}{10.9091}{mi}{105}\RG{6.775}{0.003}{10.9091}{mi}{60}\RG{18.284}{0.003}{10.9091}{mi}{106}%
}\hskip.1em\allowbreak{} ）、\allowbreak{}\hskip.1em \RB{41.133}{7.769}{3.760}{%
\RG{0.000}{0.003}{10.9091}{mi}{97}\RG{5.766}{-1.634}{7.9701}{mi}{105}\RG{8.643}{-1.634}{7.9701}{mi}{106}\RG{16.062}{0.003}{10.9091}{rm}{61}\RG{27.577}{0.003}{10.9091}{sy}{0}\RG{36.062}{0.003}{10.9091}{rm}{49}%
}\hskip.1em\allowbreak{} （\allowbreak{}\hskip.1em \RB{23.115}{7.714}{2.733}{%
\RG{0.000}{0.003}{10.9091}{mi}{105}\RG{6.775}{0.003}{10.9091}{mi}{62}\RG{18.284}{0.003}{10.9091}{mi}{106}%
}\hskip.1em\allowbreak{} ），并令\allowbreak{}\hskip.1em \RB{32.078}{7.769}{2.214}{%
\RG{0.000}{0.003}{10.9091}{mi}{97}\RG{5.767}{-1.634}{7.9701}{mi}{105}\RG{8.644}{-1.634}{7.9701}{mi}{105}\RG{15.061}{0.003}{10.9091}{rm}{61}\RG{26.559}{0.003}{10.9091}{rm}{48}%
}\hskip.1em\allowbreak{} ，则\allowbreak{}\hskip.1em \RB{8.366}{8.314}{0.500}{%
\RG{0.000}{0.003}{10.9091}{mi}{65}%
}\hskip.1em\allowbreak{} T\allowbreak{}\hskip.1em \RB{28.366}{8.314}{0.500}{%
\RG{0.000}{0.003}{10.9091}{rm}{61}\RG{11.515}{0.003}{10.9091}{sy}{0}\RG{20.000}{0.003}{10.9091}{mi}{65}%
}\hskip.1em\allowbreak{} 。因为\allowbreak{}\hskip.1em \RB{6.729}{5.320}{0.622}{%
\RG{0.000}{-0.002}{10.9091}{mi}{110}%
}\hskip.1em\allowbreak{} 为奇数，

\[
\RB{178.702}{11.431}{3.224}{%
\RT{64.852}{4.507}{7.9701}{T}\RT{0.000}{0.003}{10.9091}{d}\RT{6.066}{0.003}{10.9091}{e}\RT{10.909}{0.003}{10.9091}{t}\RG{16.971}{0.003}{10.9091}{mi}{65}\RG{28.183}{0.003}{10.9091}{rm}{61}\RT{39.698}{0.003}{10.9091}{d}\RT{45.764}{0.003}{10.9091}{e}\RT{50.608}{0.003}{10.9091}{t}\RG{56.669}{0.003}{10.9091}{mi}{65}\RG{74.166}{0.003}{10.9091}{rm}{61}\RT{85.681}{0.003}{10.9091}{d}\RT{91.746}{0.003}{10.9091}{e}\RT{96.590}{0.003}{10.9091}{t}\RG{100.834}{0.003}{10.9091}{rm}{40}\RG{105.077}{0.003}{10.9091}{sy}{0}\RG{113.562}{0.003}{10.9091}{mi}{65}\RG{121.743}{0.003}{10.9091}{rm}{41}\RG{128.998}{0.003}{10.9091}{rm}{61}\RG{140.531}{0.003}{10.9091}{sy}{0}\RT{150.834}{0.003}{10.9091}{d}\RT{156.900}{0.003}{10.9091}{e}\RT{161.743}{0.003}{10.9091}{t}\RG{167.805}{0.003}{10.9091}{mi}{65}\RG{175.987}{0.003}{10.9091}{mi}{59}%
}
\]
故det\allowbreak{}\hskip.1em \RB{28.229}{8.306}{0.745}{%
\RG{0.000}{-0.005}{10.9091}{mi}{65}\RG{11.212}{-0.005}{10.9091}{rm}{61}\RG{22.710}{-0.005}{10.9091}{rm}{48}%
}\hskip.1em\allowbreak{} 。结合已证下界，可得rank\allowbreak{}\hskip.1em \RB{56.164}{8.677}{3.232}{%
\RG{0.000}{-0.005}{10.9091}{rm}{40}\RG{4.242}{-0.005}{10.9091}{mi}{65}\RG{12.424}{-0.005}{10.9091}{rm}{41}\RG{19.679}{-0.005}{10.9091}{rm}{61}\RG{31.212}{-0.005}{10.9091}{mi}{110}\RG{40.184}{-0.005}{10.9091}{sy}{0}\RG{51.093}{-0.005}{10.9091}{rm}{49}%
}\hskip.1em\allowbreak{} 。

方法小结：模\allowbreak{}\hskip.1em \RB{5.398}{7.768}{0.500}{%
\RG{0.000}{0.002}{10.9091}{rm}{50}%
}\hskip.1em\allowbreak{} 化简消除了所有正负号的差别。先证明偶数阶矩阵可逆，再截取偶数阶主子矩阵处理奇数阶；结论还加强为“偶数阶必满秩”。

题目来源：谢启鸿《复旦大学高等代数习题集（2023版）》，第124页。

% Original page 6.
\section*{2022A10}
设\allowbreak{}\hskip.1em \RB{21.762}{8.310}{2.606}{%
\RG{0.000}{-0.001}{10.9091}{mi}{65}\RG{8.182}{-0.001}{10.9091}{mi}{59}\RG{13.015}{-0.001}{10.9091}{mi}{66}%
}\hskip.1em\allowbreak{} 为\allowbreak{}\hskip.1em \RB{6.729}{5.321}{0.621}{%
\RG{0.000}{-0.001}{10.9091}{mi}{110}%
}\hskip.1em\allowbreak{} 阶方阵，满足\allowbreak{}\hskip.1em \RB{48.730}{8.310}{0.501}{%
\RG{0.000}{-0.001}{10.9091}{mi}{65}\RG{8.182}{-0.001}{10.9091}{mi}{66}\RG{20.035}{-0.001}{10.9091}{rm}{61}\RG{31.550}{-0.001}{10.9091}{mi}{66}\RG{40.365}{-0.001}{10.9091}{mi}{65}%
}\hskip.1em\allowbreak{} 。证明：

\[
\RB{218.737}{9.293}{3.227}{%
\RT{0.000}{-0.000}{10.9091}{r}\RT{4.276}{-0.000}{10.9091}{a}\RT{9.731}{-0.000}{10.9091}{n}\RT{15.797}{-0.000}{10.9091}{k}\RG{21.557}{-0.000}{10.9091}{rm}{40}\RG{25.799}{-0.000}{10.9091}{mi}{65}\RG{33.981}{-0.000}{10.9091}{rm}{41}\RG{40.636}{-0.000}{10.9091}{rm}{43}\RT{51.557}{-0.000}{10.9091}{r}\RT{55.834}{-0.000}{10.9091}{a}\RT{61.288}{-0.000}{10.9091}{n}\RT{67.354}{-0.000}{10.9091}{k}\RG{73.114}{-0.000}{10.9091}{rm}{40}\RG{77.356}{-0.000}{10.9091}{mi}{66}\RG{86.178}{-0.000}{10.9091}{rm}{41}\RG{93.451}{-0.000}{10.9091}{sy}{21}\RT{104.966}{-0.000}{10.9091}{r}\RT{109.242}{-0.000}{10.9091}{a}\RT{114.697}{-0.000}{10.9091}{n}\RT{120.762}{-0.000}{10.9091}{k}\RG{126.523}{-0.000}{10.9091}{rm}{40}\RG{130.765}{-0.000}{10.9091}{mi}{65}\RG{142.569}{-0.000}{10.9091}{mi}{66}\RG{151.401}{-0.000}{10.9091}{rm}{41}\RG{158.056}{-0.000}{10.9091}{rm}{43}\RT{168.977}{-0.000}{10.9091}{r}\RT{173.254}{-0.000}{10.9091}{a}\RT{178.708}{-0.000}{10.9091}{n}\RT{184.774}{-0.000}{10.9091}{k}\RG{190.535}{-0.000}{10.9091}{rm}{40}\RG{194.777}{-0.000}{10.9091}{mi}{65}\RG{202.959}{-0.000}{10.9091}{mi}{66}\RG{211.781}{-0.000}{10.9091}{rm}{41}\RG{216.023}{-0.000}{10.9091}{mi}{59}%
}
\]
其中\allowbreak{}\hskip.1em \RB{28.532}{8.683}{3.226}{%
\RG{0.000}{0.001}{10.9091}{rm}{40}\RG{4.242}{0.001}{10.9091}{mi}{65}\RG{16.057}{0.001}{10.9091}{mi}{66}\RG{24.879}{0.001}{10.9091}{rm}{41}%
}\hskip.1em\allowbreak{} 表示将\allowbreak{}\hskip.1em \RB{21.762}{8.312}{2.604}{%
\RG{0.000}{0.001}{10.9091}{mi}{65}\RG{8.182}{0.001}{10.9091}{mi}{59}\RG{13.015}{0.001}{10.9091}{mi}{66}%
}\hskip.1em\allowbreak{} 水平拼接得到的\allowbreak{}\hskip.1em \RB{32.065}{7.767}{0.619}{%
\RG{0.000}{0.001}{10.9091}{mi}{110}\RG{8.972}{0.001}{10.9091}{sy}{2}\RG{19.882}{0.001}{10.9091}{rm}{50}\RG{25.336}{0.001}{10.9091}{mi}{110}%
}\hskip.1em\allowbreak{} 矩阵。

证明：设矩阵的元素属于数域\allowbreak{}\hskip.1em \RB{8.889}{7.971}{0.502}{%
\RG{0.000}{-0.002}{10.9091}{bb}{75}%
}\hskip.1em\allowbreak{} ，并令

\[
\RB{125.492}{9.292}{2.617}{%
\RG{0.000}{-0.001}{10.9091}{mi}{85}\RG{11.668}{-0.001}{10.9091}{rm}{61}\RT{23.183}{-0.001}{10.9091}{I}\RT{27.121}{-0.001}{10.9091}{m}\RG{38.026}{-0.001}{10.9091}{mi}{65}\RG{46.208}{-0.001}{10.9091}{mi}{59}\RG{72.859}{-0.001}{10.9091}{mi}{87}\RG{87.724}{-0.001}{10.9091}{rm}{61}\RT{99.239}{-0.001}{10.9091}{I}\RT{103.177}{-0.001}{10.9091}{m}\RG{114.083}{-0.001}{10.9091}{mi}{66}\RG{122.898}{-0.001}{10.9091}{mi}{58}%
}
\]
矩阵\allowbreak{}\hskip.1em \RB{28.532}{8.682}{3.227}{%
\RG{0.000}{-0.000}{10.9091}{rm}{40}\RG{4.242}{-0.000}{10.9091}{mi}{65}\RG{16.057}{-0.000}{10.9091}{mi}{66}\RG{24.879}{-0.000}{10.9091}{rm}{41}%
}\hskip.1em\allowbreak{} 的列由\allowbreak{}\hskip.1em \RB{8.366}{8.311}{0.500}{%
\RG{0.000}{-0.000}{10.9091}{mi}{65}%
}\hskip.1em\allowbreak{} 与\allowbreak{}\hskip.1em \RB{8.747}{7.951}{0.500}{%
\RG{0.000}{-0.000}{10.9091}{mi}{66}%
}\hskip.1em\allowbreak{} 的全部列组成，故

\[
\RB{243.610}{9.294}{3.226}{%
\RT{0.000}{0.001}{10.9091}{I}\RT{3.938}{0.001}{10.9091}{m}\RG{13.026}{0.001}{10.9091}{rm}{40}\RG{17.268}{0.001}{10.9091}{mi}{65}\RG{29.083}{0.001}{10.9091}{mi}{66}\RG{37.905}{0.001}{10.9091}{rm}{41}\RG{45.160}{0.001}{10.9091}{rm}{61}\RG{56.693}{0.001}{10.9091}{mi}{85}\RG{67.756}{0.001}{10.9091}{rm}{43}\RG{78.665}{0.001}{10.9091}{mi}{87}\RG{88.668}{0.001}{10.9091}{mi}{59}\RT{115.331}{0.001}{10.9091}{r}\RT{119.607}{0.001}{10.9091}{a}\RT{125.062}{0.001}{10.9091}{n}\RT{131.127}{0.001}{10.9091}{k}\RG{136.888}{0.001}{10.9091}{rm}{40}\RG{141.130}{0.001}{10.9091}{mi}{65}\RG{152.945}{0.001}{10.9091}{mi}{66}\RG{161.767}{0.001}{10.9091}{rm}{41}\RG{169.022}{0.001}{10.9091}{rm}{61}\RT{180.555}{0.001}{10.9091}{d}\RT{186.621}{0.001}{10.9091}{i}\RT{189.654}{0.001}{10.9091}{m}\RG{198.741}{0.001}{10.9091}{rm}{40}\RG{202.983}{0.001}{10.9091}{mi}{85}\RG{214.045}{0.001}{10.9091}{rm}{43}\RG{224.955}{0.001}{10.9091}{mi}{87}\RG{236.773}{0.001}{10.9091}{rm}{41}\RG{241.015}{0.001}{10.9091}{mi}{58}%
}
\]
对任意\allowbreak{}\hskip.1em \RB{33.431}{7.976}{0.934}{%
\RG{0.000}{0.003}{10.9091}{mi}{120}\RG{9.265}{0.002}{10.9091}{sy}{50}\RG{19.568}{0.003}{10.9091}{bb}{75}\RG{28.053}{3.961}{7.9701}{mi}{110}%
}\hskip.1em\allowbreak{} ，一方面\allowbreak{}\hskip.1em \RB{91.633}{8.685}{3.224}{%
\RG{0.000}{0.002}{10.9091}{mi}{65}\RG{8.182}{0.002}{10.9091}{mi}{66}\RG{16.997}{0.002}{10.9091}{mi}{120}\RG{26.269}{0.003}{10.9091}{rm}{61}\RG{37.785}{0.002}{10.9091}{mi}{65}\RG{45.966}{0.003}{10.9091}{rm}{40}\RG{50.209}{0.002}{10.9091}{mi}{66}\RG{59.024}{0.002}{10.9091}{mi}{120}\RG{65.267}{0.003}{10.9091}{rm}{41}\RG{72.539}{0.002}{10.9091}{sy}{50}\RG{82.842}{0.003}{10.9091}{mi}{85}%
}\hskip.1em\allowbreak{} ；另一方面，由\allowbreak{}\hskip.1em \RB{48.729}{8.314}{0.500}{%
\RG{0.000}{0.003}{10.9091}{mi}{65}\RG{8.182}{0.003}{10.9091}{mi}{66}\RG{20.034}{0.002}{10.9091}{rm}{61}\RG{31.549}{0.003}{10.9091}{mi}{66}\RG{40.364}{0.003}{10.9091}{mi}{65}%
}\hskip.1em\allowbreak{} ，

\[
\RB{133.224}{8.686}{3.223}{%
\RG{0.000}{0.004}{10.9091}{mi}{65}\RG{8.182}{0.004}{10.9091}{mi}{66}\RG{16.997}{0.004}{10.9091}{mi}{120}\RG{26.269}{0.004}{10.9091}{rm}{61}\RG{37.784}{0.004}{10.9091}{mi}{66}\RG{46.599}{0.004}{10.9091}{mi}{65}\RG{54.781}{0.004}{10.9091}{mi}{120}\RG{64.054}{0.004}{10.9091}{rm}{61}\RG{75.569}{0.004}{10.9091}{mi}{66}\RG{84.391}{0.004}{10.9091}{rm}{40}\RG{88.634}{0.004}{10.9091}{mi}{65}\RG{96.816}{0.004}{10.9091}{mi}{120}\RG{103.051}{0.004}{10.9091}{rm}{41}\RG{110.323}{0.004}{10.9091}{sy}{50}\RG{120.626}{0.004}{10.9091}{mi}{87}\RG{130.630}{0.004}{10.9091}{mi}{58}%
}
\]
因此

\[
\RB{233.920}{9.293}{3.227}{%
\RT{0.000}{-0.000}{10.9091}{I}\RT{3.938}{-0.000}{10.9091}{m}\RG{13.026}{-0.000}{10.9091}{rm}{40}\RG{17.268}{-0.000}{10.9091}{mi}{65}\RG{25.450}{-0.000}{10.9091}{mi}{66}\RG{34.272}{-0.000}{10.9091}{rm}{41}\RG{41.545}{-0.000}{10.9091}{sy}{18}\RG{53.060}{-0.000}{10.9091}{mi}{85}\RG{64.123}{-0.000}{10.9091}{sy}{92}\RG{73.820}{-0.000}{10.9091}{mi}{87}\RG{83.823}{-0.000}{10.9091}{mi}{59}\RT{110.486}{-0.000}{10.9091}{r}\RT{114.762}{-0.000}{10.9091}{a}\RT{120.217}{-0.000}{10.9091}{n}\RT{126.282}{-0.000}{10.9091}{k}\RG{132.043}{-0.000}{10.9091}{rm}{40}\RG{136.285}{-0.000}{10.9091}{mi}{65}\RG{144.467}{-0.000}{10.9091}{mi}{66}\RG{153.289}{-0.000}{10.9091}{rm}{41}\RG{160.562}{-0.000}{10.9091}{sy}{20}\RT{172.077}{-0.000}{10.9091}{d}\RT{178.143}{-0.000}{10.9091}{i}\RT{181.176}{-0.000}{10.9091}{m}\RG{190.263}{-0.000}{10.9091}{rm}{40}\RG{194.506}{-0.000}{10.9091}{mi}{85}\RG{205.568}{-0.000}{10.9091}{sy}{92}\RG{215.265}{-0.000}{10.9091}{mi}{87}\RG{227.083}{-0.000}{10.9091}{rm}{41}\RG{231.325}{-0.000}{10.9091}{mi}{58}%
}
\]
利用子空间的维数公式，得到

\[
\RB{258.054}{9.296}{39.553}{%
\RT{0.000}{0.003}{10.9091}{r}\RT{4.276}{0.003}{10.9091}{a}\RT{9.731}{0.003}{10.9091}{n}\RT{15.797}{0.003}{10.9091}{k}\RG{21.556}{0.003}{10.9091}{rm}{40}\RG{25.799}{0.003}{10.9091}{mi}{65}\RG{37.603}{0.003}{10.9091}{mi}{66}\RG{46.435}{0.003}{10.9091}{rm}{41}\RG{53.090}{0.003}{10.9091}{rm}{43}\RT{64.012}{0.003}{10.9091}{r}\RT{68.288}{0.003}{10.9091}{a}\RT{73.743}{0.003}{10.9091}{n}\RT{79.808}{0.003}{10.9091}{k}\RG{85.568}{0.003}{10.9091}{rm}{40}\RG{89.811}{0.003}{10.9091}{mi}{65}\RG{97.993}{0.003}{10.9091}{mi}{66}\RG{106.815}{0.003}{10.9091}{rm}{41}\RG{114.087}{0.003}{10.9091}{sy}{20}\RT{125.602}{0.003}{10.9091}{d}\RT{131.668}{0.003}{10.9091}{i}\RT{134.700}{0.003}{10.9091}{m}\RG{143.788}{0.003}{10.9091}{rm}{40}\RG{148.030}{0.003}{10.9091}{mi}{85}\RG{159.092}{0.003}{10.9091}{rm}{43}\RG{170.001}{0.003}{10.9091}{mi}{87}\RG{181.821}{0.003}{10.9091}{rm}{41}\RG{188.475}{0.003}{10.9091}{rm}{43}\RT{199.396}{0.003}{10.9091}{d}\RT{205.461}{0.003}{10.9091}{i}\RT{208.494}{0.003}{10.9091}{m}\RG{217.582}{0.003}{10.9091}{rm}{40}\RG{221.824}{0.003}{10.9091}{mi}{85}\RG{232.886}{0.003}{10.9091}{sy}{92}\RG{242.583}{0.003}{10.9091}{mi}{87}\RG{254.401}{0.003}{10.9091}{rm}{41}\RG{114.087}{-18.161}{10.9091}{rm}{61}\RT{125.602}{-18.161}{10.9091}{d}\RT{131.668}{-18.161}{10.9091}{i}\RT{134.700}{-18.161}{10.9091}{m}\RG{145.606}{-18.161}{10.9091}{mi}{85}\RG{156.668}{-18.161}{10.9091}{rm}{43}\RT{167.577}{-18.161}{10.9091}{d}\RT{173.643}{-18.161}{10.9091}{i}\RT{176.676}{-18.161}{10.9091}{m}\RG{187.582}{-18.161}{10.9091}{mi}{87}\RG{114.087}{-36.325}{10.9091}{rm}{61}\RT{125.602}{-36.325}{10.9091}{r}\RT{129.878}{-36.325}{10.9091}{a}\RT{135.333}{-36.325}{10.9091}{n}\RT{141.399}{-36.325}{10.9091}{k}\RG{147.159}{-36.325}{10.9091}{rm}{40}\RG{151.401}{-36.325}{10.9091}{mi}{65}\RG{159.583}{-36.325}{10.9091}{rm}{41}\RG{166.238}{-36.325}{10.9091}{rm}{43}\RT{177.160}{-36.325}{10.9091}{r}\RT{181.436}{-36.325}{10.9091}{a}\RT{186.891}{-36.325}{10.9091}{n}\RT{192.956}{-36.325}{10.9091}{k}\RG{198.716}{-36.325}{10.9091}{rm}{40}\RG{202.959}{-36.325}{10.9091}{mi}{66}\RG{211.781}{-36.325}{10.9091}{rm}{41}\RG{216.023}{-36.325}{10.9091}{mi}{58}%
}
\]
所用维数公式可由基的扩充直接证明：先取\allowbreak{}\hskip.1em \RB{32.237}{7.947}{0.744}{%
\RG{0.000}{-0.004}{10.9091}{mi}{85}\RG{10.835}{-0.004}{10.9091}{sy}{92}\RG{20.304}{-0.004}{10.9091}{mi}{87}%
}\hskip.1em\allowbreak{} 的一组基，再分别扩充成\allowbreak{}\hskip.1em \RB{8.791}{7.947}{0.744}{%
\RG{0.000}{-0.004}{10.9091}{mi}{85}%
}\hskip.1em\allowbreak{} 与\allowbreak{}\hskip.1em \RB{11.933}{7.947}{0.744}{%
\RG{0.000}{-0.004}{10.9091}{mi}{87}%
}\hskip.1em\allowbreak{} 的基；把这两组基中公共部分只保留一次，所得向量组就是\allowbreak{}\hskip.1em \RB{33.904}{7.952}{1.404}{%
\RG{0.000}{0.001}{10.9091}{mi}{85}\RG{11.062}{0.001}{10.9091}{rm}{43}\RG{21.971}{0.001}{10.9091}{mi}{87}%
}\hskip.1em\allowbreak{} 的一组基。\allowbreak{}\hskip.1em \RB{8.366}{7.996}{0.500}{%
\RG{0.000}{0.001}{10.9091}{am}{3}%
}\hskip.1em\allowbreak{}补充：上述证明还表明，等号成立当且仅当

\[
\RB{114.460}{9.290}{3.230}{%
\RT{0.000}{-0.003}{10.9091}{I}\RT{3.938}{-0.003}{10.9091}{m}\RG{13.025}{-0.003}{10.9091}{rm}{40}\RG{17.268}{-0.003}{10.9091}{mi}{65}\RG{25.450}{-0.003}{10.9091}{mi}{66}\RG{34.272}{-0.003}{10.9091}{rm}{41}\RG{41.527}{-0.003}{10.9091}{rm}{61}\RT{53.059}{-0.003}{10.9091}{I}\RT{56.998}{-0.003}{10.9091}{m}\RG{67.904}{-0.003}{10.9091}{mi}{65}\RG{78.510}{-0.003}{10.9091}{sy}{92}\RT{88.207}{-0.003}{10.9091}{I}\RT{92.145}{-0.003}{10.9091}{m}\RG{103.051}{-0.003}{10.9091}{mi}{66}\RG{111.865}{-0.003}{10.9091}{mi}{58}%
}
\]
这是因为证明中唯一的不等式来自Im\allowbreak{}\hskip.1em \RB{73.538}{8.680}{3.229}{%
\RG{0.000}{-0.002}{10.9091}{rm}{40}\RG{4.242}{-0.002}{10.9091}{mi}{65}\RG{12.424}{-0.002}{10.9091}{mi}{66}\RG{21.246}{-0.002}{10.9091}{rm}{41}\RG{28.809}{-0.002}{10.9091}{sy}{18}\RG{40.614}{-0.002}{10.9091}{mi}{85}\RG{51.792}{-0.002}{10.9091}{sy}{92}\RG{61.606}{-0.002}{10.9091}{mi}{87}%
}\hskip.1em\allowbreak{} ，有限维子空间在包含关系下维数相等当且仅当两者相等。

方法小结：将水平拼接矩阵的像解释为两个像空间之和，将乘积矩阵的像放进它们的交。交换条件恰好用于第二个包含关系，之后只需交与和的维数公式。

题目来源：谢启鸿《复旦大学高等代数习题集（2023版）》，第124页。

% Original page 7.
\section*{2022A11}
设\allowbreak{}\hskip.1em \RB{35.464}{8.310}{2.606}{%
\RG{0.000}{-0.001}{10.9091}{mi}{65}\RG{8.182}{-0.001}{10.9091}{mi}{59}\RG{13.015}{-0.001}{10.9091}{mi}{66}\RG{21.829}{-0.001}{10.9091}{mi}{59}\RG{26.673}{-0.001}{10.9091}{mi}{67}%
}\hskip.1em\allowbreak{} 分别为数域\allowbreak{}\hskip.1em \RB{8.889}{7.972}{0.501}{%
\RG{0.000}{-0.001}{10.9091}{bb}{75}%
}\hskip.1em\allowbreak{} 上的\allowbreak{}\hskip.1em \RB{29.641}{5.855}{0.621}{%
\RG{0.000}{-0.001}{10.9091}{mi}{109}\RG{12.003}{-0.001}{10.9091}{sy}{2}\RG{22.912}{-0.001}{10.9091}{mi}{110}%
}\hskip.1em\allowbreak{} 、\allowbreak{}\hskip.1em \RB{24.253}{5.855}{2.617}{%
\RG{0.000}{-0.001}{10.9091}{mi}{112}\RG{7.913}{-0.001}{10.9091}{sy}{2}\RG{18.822}{-0.001}{10.9091}{mi}{113}%
}\hskip.1em\allowbreak{} 、\allowbreak{}\hskip.1em \RB{28.343}{5.855}{2.617}{%
\RG{0.000}{-0.001}{10.9091}{mi}{109}\RG{12.003}{-0.001}{10.9091}{sy}{2}\RG{22.912}{-0.001}{10.9091}{mi}{113}%
}\hskip.1em\allowbreak{} 矩阵，且矩阵方程

\[
\RB{106.509}{8.311}{2.605}{%
\RG{0.000}{-0.000}{10.9091}{mi}{65}\RG{8.182}{-0.000}{10.9091}{mi}{89}\RG{19.969}{-0.000}{10.9091}{rm}{61}\RG{31.484}{-0.000}{10.9091}{mi}{67}\RG{39.448}{-0.000}{10.9091}{mi}{59}\RG{66.099}{-0.000}{10.9091}{mi}{90}\RG{74.325}{-0.000}{10.9091}{mi}{66}\RG{86.203}{-0.000}{10.9091}{rm}{61}\RG{97.718}{-0.000}{10.9091}{mi}{67}%
}
\]
都有解：证明：矩阵方程\allowbreak{}\hskip.1em \RB{50.234}{8.312}{0.739}{%
\RG{0.000}{0.001}{10.9091}{mi}{65}\RG{8.182}{0.001}{10.9091}{mi}{88}\RG{18.066}{0.001}{10.9091}{mi}{66}\RG{29.928}{0.001}{10.9091}{rm}{61}\RG{41.443}{0.001}{10.9091}{mi}{67}%
}\hskip.1em\allowbreak{} 也有解，其中\allowbreak{}\hskip.1em \RB{34.275}{7.952}{2.604}{%
\RG{0.000}{0.001}{10.9091}{mi}{88}\RG{9.273}{0.001}{10.9091}{mi}{59}\RG{14.116}{0.001}{10.9091}{mi}{89}\RG{21.055}{0.001}{10.9091}{mi}{59}\RG{25.888}{0.001}{10.9091}{mi}{90}%
}\hskip.1em\allowbreak{} 的大小分别为\allowbreak{}\hskip.1em \RB{25.728}{5.857}{2.615}{%
\RG{0.000}{0.001}{10.9091}{mi}{110}\RG{8.973}{0.001}{10.9091}{sy}{2}\RG{19.882}{0.001}{10.9091}{mi}{112}%
}\hskip.1em\allowbreak{} 、\allowbreak{}\hskip.1em \RB{25.313}{5.857}{2.615}{%
\RG{0.000}{0.001}{10.9091}{mi}{110}\RG{8.973}{0.001}{10.9091}{sy}{2}\RG{19.882}{0.001}{10.9091}{mi}{113}%
}\hskip.1em\allowbreak{} 、\allowbreak{}\hskip.1em \RB{28.758}{5.857}{2.615}{%
\RG{0.000}{0.001}{10.9091}{mi}{109}\RG{12.003}{0.001}{10.9091}{sy}{2}\RG{22.912}{0.001}{10.9091}{mi}{112}%
}\hskip.1em\allowbreak{} 。

证明：取已知方程的两个解\allowbreak{}\hskip.1em \RB{27.759}{7.949}{2.607}{%
\RG{0.000}{-0.002}{10.9091}{mi}{89}\RG{6.333}{-1.638}{7.9701}{rm}{48}\RG{11.065}{-0.002}{10.9091}{mi}{59}\RG{15.898}{-0.002}{10.9091}{mi}{90}\RG{23.361}{-1.638}{7.9701}{rm}{48}%
}\hskip.1em\allowbreak{} ，使

\[
\RB{114.526}{8.310}{2.606}{%
\RG{0.000}{-0.001}{10.9091}{mi}{65}\RG{8.182}{-0.001}{10.9091}{mi}{89}\RG{22.278}{-0.001}{10.9091}{rm}{61}\RG{33.793}{-0.001}{10.9091}{mi}{67}\RG{41.757}{-0.001}{10.9091}{mi}{59}\RG{68.408}{-0.001}{10.9091}{mi}{90}\RG{80.611}{-0.001}{10.9091}{mi}{66}\RG{92.464}{-0.001}{10.9091}{rm}{61}\RG{103.979}{-0.001}{10.9091}{mi}{67}\RG{111.932}{-0.001}{10.9091}{mi}{58}\RG{14.516}{-1.637}{7.9701}{rm}{48}\RG{75.879}{-1.637}{7.9701}{rm}{48}%
}
\]
关键是构造一个\allowbreak{}\hskip.1em \RB{28.346}{5.856}{2.616}{%
\RG{0.000}{-0.000}{10.9091}{mi}{113}\RG{7.686}{-0.000}{10.9091}{sy}{2}\RG{18.595}{-0.000}{10.9091}{mi}{109}%
}\hskip.1em\allowbreak{} 矩阵\allowbreak{}\hskip.1em \RB{8.180}{7.886}{0.500}{%
\RG{0.000}{-0.000}{10.9091}{mi}{84}%
}\hskip.1em\allowbreak{} ，满足

\[
\RB{50.138}{8.186}{0.745}{%
\RG{0.000}{-0.005}{10.9091}{mi}{67}\RG{8.564}{-0.005}{10.9091}{mi}{84}\RG{16.451}{-0.005}{10.9091}{mi}{67}\RG{28.075}{-0.005}{10.9091}{rm}{61}\RG{39.590}{-0.005}{10.9091}{mi}{67}\RG{47.543}{-0.005}{10.9091}{mi}{58}%
}
\]
一旦找到这样的\allowbreak{}\hskip.1em \RB{8.180}{7.884}{0.502}{%
\RG{0.000}{-0.002}{10.9091}{mi}{84}%
}\hskip.1em\allowbreak{} ，令

\[
\RB{62.166}{7.954}{2.602}{%
\RG{0.000}{0.003}{10.9091}{mi}{88}\RG{16.801}{0.003}{10.9091}{rm}{61}\RG{28.316}{0.003}{10.9091}{mi}{89}\RG{39.382}{0.003}{10.9091}{mi}{84}\RG{47.259}{0.003}{10.9091}{mi}{90}\RG{59.451}{0.003}{10.9091}{mi}{59}\RG{9.038}{-1.634}{7.9701}{rm}{48}\RG{34.649}{-1.634}{7.9701}{rm}{48}\RG{54.719}{-1.634}{7.9701}{rm}{48}%
}
\]
便有

\[
\RB{175.233}{8.682}{3.227}{%
\RG{0.000}{-0.000}{10.9091}{mi}{65}\RG{8.182}{-0.000}{10.9091}{mi}{88}\RG{21.952}{-0.000}{10.9091}{mi}{66}\RG{33.805}{-0.000}{10.9091}{rm}{61}\RG{45.304}{-0.000}{10.9091}{rm}{40}\RG{49.562}{-0.000}{10.9091}{mi}{65}\RG{57.744}{-0.000}{10.9091}{mi}{89}\RG{68.811}{-0.000}{10.9091}{rm}{41}\RG{73.053}{-0.000}{10.9091}{mi}{84}\RG{80.943}{-0.000}{10.9091}{rm}{40}\RG{85.186}{-0.000}{10.9091}{mi}{90}\RG{97.365}{-0.000}{10.9091}{mi}{66}\RG{106.187}{-0.000}{10.9091}{rm}{41}\RG{113.441}{-0.000}{10.9091}{rm}{61}\RG{124.975}{-0.000}{10.9091}{mi}{67}\RG{133.539}{-0.000}{10.9091}{mi}{84}\RG{141.426}{-0.000}{10.9091}{mi}{67}\RG{153.050}{-0.000}{10.9091}{rm}{61}\RG{164.565}{-0.000}{10.9091}{mi}{67}\RG{172.518}{-0.000}{10.9091}{mi}{59}\RG{17.220}{-1.637}{7.9701}{rm}{48}\RG{64.078}{-1.637}{7.9701}{rm}{48}\RG{92.633}{-1.637}{7.9701}{rm}{48}%
}
\]
而\allowbreak{}\hskip.1em \RB{30.801}{7.953}{2.301}{%
\RG{0.000}{0.002}{10.9091}{mi}{89}\RG{6.333}{-1.634}{7.9701}{rm}{48}\RG{11.065}{0.002}{10.9091}{mi}{84}\RG{18.942}{0.002}{10.9091}{mi}{90}\RG{26.403}{-1.634}{7.9701}{rm}{48}%
}\hskip.1em\allowbreak{} 的大小正是\allowbreak{}\hskip.1em \RB{25.727}{5.858}{2.614}{%
\RG{0.000}{0.002}{10.9091}{mi}{110}\RG{8.972}{0.002}{10.9091}{sy}{2}\RG{19.881}{0.002}{10.9091}{mi}{112}%
}\hskip.1em\allowbreak{} 。

下面证明所需的\allowbreak{}\hskip.1em \RB{8.180}{7.883}{0.503}{%
\RG{0.000}{-0.003}{10.9091}{mi}{84}%
}\hskip.1em\allowbreak{} 总是存在。若\allowbreak{}\hskip.1em \RB{28.624}{8.188}{0.743}{%
\RG{0.000}{-0.003}{10.9091}{mi}{67}\RG{11.607}{-0.003}{10.9091}{rm}{61}\RG{23.105}{-0.003}{10.9091}{rm}{48}%
}\hskip.1em\allowbreak{} ，直接取\allowbreak{}\hskip.1em \RB{33.818}{7.948}{2.306}{%
\RG{0.000}{-0.003}{10.9091}{mi}{88}\RG{9.038}{-1.639}{7.9701}{rm}{48}\RG{16.801}{-0.003}{10.9091}{rm}{61}\RG{28.299}{-0.003}{10.9091}{rm}{48}%
}\hskip.1em\allowbreak{} 即可。以下设\allowbreak{}\hskip.1em \RB{16.621}{5.319}{0.623}{%
\RG{0.000}{-0.003}{10.9091}{mi}{114}\RG{8.255}{-0.003}{10.9091}{rm}{61}%
}\hskip.1em\allowbreak{} rank\allowbreak{}\hskip.1em \RB{37.126}{8.679}{3.230}{%
\RG{0.000}{-0.003}{10.9091}{rm}{40}\RG{4.242}{-0.003}{10.9091}{mi}{67}\RG{12.820}{-0.003}{10.9091}{rm}{41}\RG{20.092}{-0.003}{10.9091}{mi}{62}\RG{31.607}{-0.003}{10.9091}{rm}{48}%
}\hskip.1em\allowbreak{} 。对\allowbreak{}\hskip.1em \RB{8.791}{8.193}{0.738}{%
\RG{0.000}{0.002}{10.9091}{mi}{67}%
}\hskip.1em\allowbreak{} 施行可逆的初等行、列变换，可将它化为秩标准形，即存在可逆矩阵\allowbreak{}\hskip.1em \RB{52.143}{8.502}{0.934}{%
\RG{0.000}{0.002}{10.9091}{mi}{80}\RG{11.549}{0.002}{10.9091}{sy}{50}\RG{21.852}{0.002}{10.9091}{bb}{75}\RG{30.337}{3.961}{7.9701}{mi}{109}\RG{37.827}{3.961}{7.9701}{sy}{2}\RG{44.414}{3.961}{7.9701}{mi}{109}%
}\hskip.1em\allowbreak{} 、\allowbreak{}\hskip.1em \RB{45.405}{8.502}{2.614}{%
\RG{0.000}{0.002}{10.9091}{mi}{81}\RG{11.654}{0.002}{10.9091}{sy}{50}\RG{21.957}{0.002}{10.9091}{bb}{75}\RG{30.442}{3.961}{7.9701}{mi}{113}\RG{34.509}{3.961}{7.9701}{sy}{2}\RG{41.095}{3.961}{7.9701}{mi}{113}%
}\hskip.1em\allowbreak{} ，使

\[
\RB{129.242}{19.591}{16.181}{%
\RG{64.147}{18.654}{10.9092}{ex}{32}\RG{97.550}{18.654}{10.9092}{ex}{33}\RG{72.783}{7.279}{10.9091}{mi}{73}\RG{92.096}{7.279}{10.9091}{rm}{48}\RG{77.579}{5.643}{7.9701}{mi}{114}\RG{0.000}{-0.001}{10.9091}{mi}{80}\RG{8.509}{-0.001}{10.9091}{mi}{67}\RG{17.084}{-0.001}{10.9091}{mi}{81}\RG{28.750}{-0.001}{10.9091}{rm}{61}\RG{40.265}{-0.001}{10.9091}{mi}{68}\RG{52.631}{-0.001}{10.9091}{rm}{61}\RG{126.647}{-0.001}{10.9091}{mi}{58}\RG{74.731}{-7.896}{10.9091}{rm}{48}\RG{92.087}{-7.896}{10.9091}{rm}{48}\RG{106.187}{-14.135}{7.9701}{mi}{109}\RG{113.677}{-14.135}{7.9701}{sy}{2}\RG{120.264}{-14.135}{7.9701}{mi}{113}%
}
\]
具体地，高斯消元先选出\allowbreak{}\hskip.1em \RB{5.256}{5.317}{0.625}{%
\RG{0.000}{-0.005}{10.9091}{mi}{114}%
}\hskip.1em\allowbreak{} 个主元，通过换行、换列将它们置于前\allowbreak{}\hskip.1em \RB{5.256}{5.317}{0.625}{%
\RG{0.000}{-0.005}{10.9091}{mi}{114}%
}\hskip.1em\allowbreak{} 个位置，再将主元化为\allowbreak{}\hskip.1em \RB{5.071}{7.761}{0.505}{%
\RG{0.000}{-0.005}{10.9091}{rm}{49}%
}\hskip.1em\allowbreak{} 并消去相应行、列的其余元素，就得到这个标准形。

令\allowbreak{}\hskip.1em \RB{32.498}{7.946}{0.505}{%
\RG{0.000}{-0.005}{10.9091}{mi}{69}\RG{11.712}{-0.005}{10.9091}{rm}{61}\RG{23.227}{-0.005}{10.9091}{mi}{68}%
}\hskip.1em\allowbreak{} T\allowbreak{}\hskip.1em \RB{37.170}{8.495}{0.941}{%
\RG{0.000}{-0.005}{10.9091}{sy}{50}\RG{10.303}{-0.005}{10.9091}{bb}{75}\RG{18.788}{3.954}{7.9701}{mi}{113}\RG{22.854}{3.954}{7.9701}{sy}{2}\RG{29.441}{3.954}{7.9701}{mi}{109}%
}\hskip.1em\allowbreak{} 。由\allowbreak{}\hskip.1em \RB{9.271}{7.946}{0.505}{%
\RG{0.000}{-0.005}{10.9091}{mi}{68}%
}\hskip.1em\allowbreak{} 的形状，直接相乘得

\[
\RB{53.808}{7.947}{0.504}{%
\RG{0.000}{-0.004}{10.9091}{mi}{68}\RG{9.316}{-0.004}{10.9091}{mi}{69}\RG{18.000}{-0.004}{10.9091}{mi}{68}\RG{30.382}{-0.004}{10.9091}{rm}{61}\RG{41.897}{-0.004}{10.9091}{mi}{68}\RG{51.214}{-0.004}{10.9091}{mi}{58}%
}
\]
现在取\allowbreak{}\hskip.1em \RB{48.453}{8.188}{2.619}{%
\RG{0.000}{-0.003}{10.9091}{mi}{84}\RG{10.920}{-0.003}{10.9091}{rm}{61}\RG{22.436}{-0.003}{10.9091}{mi}{81}\RG{31.054}{-0.003}{10.9091}{mi}{69}\RG{39.727}{-0.003}{10.9091}{mi}{80}%
}\hskip.1em\allowbreak{} 。由于\allowbreak{}\hskip.1em \RB{71.569}{9.756}{2.619}{%
\RG{0.000}{-0.003}{10.9091}{mi}{67}\RG{11.608}{-0.003}{10.9091}{rm}{61}\RG{23.123}{-0.003}{10.9091}{mi}{80}\RG{31.642}{3.956}{7.9701}{sy}{0}\RG{38.229}{3.956}{7.9701}{rm}{49}\RG{42.961}{-0.003}{10.9091}{mi}{68}\RG{52.278}{-0.003}{10.9091}{mi}{81}\RG{60.921}{3.956}{7.9701}{sy}{0}\RG{67.507}{3.956}{7.9701}{rm}{49}%
}\hskip.1em\allowbreak{} ，有

\[
\RB{188.513}{10.302}{38.947}{%
\RG{52.352}{4.502}{7.9701}{sy}{0}\RG{58.938}{4.502}{7.9701}{rm}{49}\RG{81.631}{4.502}{7.9701}{sy}{0}\RG{88.217}{4.502}{7.9701}{rm}{49}\RG{144.263}{4.502}{7.9701}{sy}{0}\RG{150.850}{4.502}{7.9701}{rm}{49}\RG{173.541}{4.502}{7.9701}{sy}{0}\RG{180.128}{4.502}{7.9701}{rm}{49}\RG{0.000}{-0.002}{10.9091}{mi}{67}\RG{8.564}{-0.002}{10.9091}{mi}{84}\RG{16.451}{-0.002}{10.9091}{mi}{67}\RG{28.075}{-0.002}{10.9091}{rm}{61}\RG{39.574}{-0.002}{10.9091}{rm}{40}\RG{43.833}{-0.002}{10.9091}{mi}{80}\RG{63.672}{-0.002}{10.9091}{mi}{68}\RG{72.988}{-0.002}{10.9091}{mi}{81}\RG{92.950}{-0.002}{10.9091}{rm}{41}\RG{97.183}{-0.002}{10.9091}{rm}{40}\RG{101.434}{-0.002}{10.9091}{mi}{81}\RG{110.052}{-0.002}{10.9091}{mi}{69}\RG{118.725}{-0.002}{10.9091}{mi}{80}\RG{127.259}{-0.002}{10.9091}{rm}{41}\RG{131.492}{-0.002}{10.9091}{rm}{40}\RG{135.744}{-0.002}{10.9091}{mi}{80}\RG{155.582}{-0.002}{10.9091}{mi}{68}\RG{164.899}{-0.002}{10.9091}{mi}{81}\RG{184.860}{-0.002}{10.9091}{rm}{41}\RG{48.109}{-13.662}{7.9701}{sy}{0}\RG{54.696}{-13.662}{7.9701}{rm}{49}\RG{95.405}{-13.662}{7.9701}{sy}{0}\RG{101.991}{-13.662}{7.9701}{rm}{49}\RG{28.075}{-18.166}{10.9091}{rm}{61}\RG{39.590}{-18.166}{10.9091}{mi}{80}\RG{59.428}{-18.166}{10.9091}{mi}{68}\RG{68.745}{-18.166}{10.9091}{mi}{69}\RG{77.429}{-18.166}{10.9091}{mi}{68}\RG{86.756}{-18.166}{10.9091}{mi}{81}\RG{48.109}{-31.825}{7.9701}{sy}{0}\RG{54.696}{-31.825}{7.9701}{rm}{49}\RG{77.388}{-31.825}{7.9701}{sy}{0}\RG{83.975}{-31.825}{7.9701}{rm}{49}\RG{28.075}{-36.330}{10.9091}{rm}{61}\RG{39.590}{-36.330}{10.9091}{mi}{80}\RG{59.428}{-36.330}{10.9091}{mi}{68}\RG{68.745}{-36.330}{10.9091}{mi}{81}\RG{91.737}{-36.330}{10.9091}{rm}{61}\RG{103.253}{-36.330}{10.9091}{mi}{67}\RG{111.206}{-36.330}{10.9091}{mi}{58}%
}
\]
所需\allowbreak{}\hskip.1em \RB{8.180}{7.887}{0.500}{%
\RG{0.000}{0.001}{10.9091}{mi}{84}%
}\hskip.1em\allowbreak{} 已构造出来，故\allowbreak{}\hskip.1em \RB{77.051}{8.192}{2.615}{%
\RG{0.000}{0.001}{10.9091}{mi}{88}\RG{9.038}{-1.635}{7.9701}{rm}{48}\RG{16.801}{0.001}{10.9091}{rm}{61}\RG{28.316}{0.001}{10.9091}{mi}{89}\RG{34.649}{-1.635}{7.9701}{rm}{48}\RG{39.381}{0.001}{10.9091}{mi}{81}\RG{48.000}{0.001}{10.9091}{mi}{69}\RG{56.672}{0.001}{10.9091}{mi}{80}\RG{65.192}{0.001}{10.9091}{mi}{90}\RG{72.654}{-1.635}{7.9701}{rm}{48}%
}\hskip.1em\allowbreak{} 为原方程的一个解：\allowbreak{}\hskip.1em \RB{8.366}{7.996}{0.500}{%
\RG{0.000}{0.001}{10.9091}{am}{3}%
}\hskip.1em\allowbreak{}

方法小结：先用秩标准形构造满足\allowbreak{}\hskip.1em \RB{48.871}{8.190}{0.741}{%
\RG{0.000}{-0.001}{10.9091}{mi}{67}\RG{8.564}{-0.001}{10.9091}{mi}{84}\RG{16.451}{-0.001}{10.9091}{mi}{67}\RG{28.320}{-0.001}{10.9091}{rm}{61}\RG{40.080}{-0.001}{10.9091}{mi}{67}%
}\hskip.1em\allowbreak{} 的矩阵，再把两个单边方程的解组合为\allowbreak{}\hskip.1em \RB{30.800}{7.950}{2.305}{%
\RG{0.000}{-0.001}{10.9091}{mi}{89}\RG{6.333}{-1.638}{7.9701}{rm}{48}\RG{11.065}{-0.001}{10.9091}{mi}{84}\RG{18.942}{-0.001}{10.9091}{mi}{90}\RG{26.402}{-1.638}{7.9701}{rm}{48}%
}\hskip.1em\allowbreak{} 。这样既证明了双边方程有解，也提供了求解的具体办法。

题目来源：谢启鸿《复旦大学高等代数习题集（2023版）》，第124页。


\end{document}
